%% Theory of Computing
%% File: v016a005.tex
%% Publication date: September 25, 2020
%% Authors: Alessandro Chiesa, Peter Manohar, and Igor Shinkar

%

%
%
%
%

%

%

%
%
%

%
%
%

%
%
%
%
%
%
%
%
%
%
%
%
%
%
%
%
%
%
%
%
%
%
%
%

%

\documentclass{toc}

%
%
%
\tocdetails{%
  title = {On Axis-Parallel Tests for Tensor Product Codes},
%
%
%
%
  number_of_pages = {29},
  number_of_bibitems = {36},
  number_of_figures = {0},
  conference_version = {RANDOM17}, %
    %
    %
%
%
%
  author = {Alessandro Chiesa, Peter Manohar, and Igor Shinkar},
    %
    %
    %
  authorlist = {Alessandro Chiesa, Peter Manohar, Igor Shinkar},
    %
    %
    %
    %
    %
    %
    %
%
%
  %
  %
  %
  %
%
  %
  acmclassification = {E.4, F.2.2},             %
        %
    %
    %
  %
  amsclassification = {68P30, 94B05, 94B25},                %
    %
    %
  %
  keywords = {error-correcting codes, tensor product codes, locally testable codes, low-degree testing, extremal graph theory}, %
  %
  %
    %
    %
    %
    %
    %
    %
}   %

    %
    %
    %
    %
    %
    %
    %
    %
    %
    %
    %
    %
    %
    %
    %
    %
    %
    %
    %
    %
    %
    %
    %
    %
    %
    %
    %
    %
    %

%
%
\tocdetails{%
volume = 16,
number = 5,
year = 2020,
%
%
%
%
%
%
%
received = {August 30, 2018},
revised = {July 22, 2020},
published = {September 25, 2020},
%
%
doi = {10.4086/toc.2020.v016a005},
%
}   %

\usepackage{paralist}

\theoremstyle{plain}
\newtheorem{ftheorem}[theorem]{Theorem}

\begin{document}


\newcommand{\colored}{}
\newcommand{\reldist}{\colored{\delta}}
\newcommand{\robust}{\colored{\rho}} %
\newcommand{\tester}{\colored{\mathcal H}} %
\newcommand{\plane}{{\colored{H}}} %
\newcommand{\otherplane}{\colored{ {H'}}} %
\newcommand{\hamdist}{\colored{\Delta}} %
\newcommand{\codedist}{\colored{d}} %
\newcommand{\blocklength}{\colored{n}} %
\newcommand{\testerinput}{\colored{M}} %
\newcommand{\closestlocalword}[1]{ \colored{g_{#1}}} %
\newcommand{\closestglobalword}{\colored{f}} %
\newcommand{\noteq}{\colored{\text{neq}}}
\newcommand{\arbtester}{\colored{\mathcal T}}
\newcommand{\Field}{\colored{\mathbb F}}
\newcommand{\st}{\colored{\text{ s.t. }}}
\newcommand{\word}{\colored{w}}
\newcommand{\otherword}{\colored{v}}
%

%
\newcommand{\degree}{\colored{r}}
\newcommand{\rowpoly}{\colored{\mathcal R}}
\newcommand{\colpoly}{\colored{\mathcal C}}
\newcommand{\prox}{\colored{\delta}}
\newcommand{\PSprox}{\colored{\gamma}}
\newcommand{\numrect}{\colored{u}}

\newcommand{\DefineEqual}{:=}
\newcommand{\floor}[1]{{\lfloor#1\rfloor}}
\newcommand{\ceil}[1]{{\lceil#1\rceil}}
\newcommand{\poly}{\operatorname{poly}}
\newcommand{\polylog}{\operatorname{polylog}}
\newcommand{\E}{{\mathbb E}}
\newcommand{\eps}{\varepsilon}
\newcommand\abs[1]{\colored{\left|#1\right|}}

\newcommand{\Naturals}{\mathbb{N}}

\newcommand{\parhead}{\paragraph}

\begin{frontmatter} %

%
  \title{On Axis-Parallel Tests for \\ Tensor Product Codes\titlefootnote{A preliminary version of this paper appeared in the
\href{https://doi.org/10.4230/LIPIcs.APPROX-RANDOM.2017.39}{Proceedings
 of the 21st International Workshop on Randomization and Computation 
(RANDOM'17).}}} 
  %

%
%

%
%
%

%
%
%
%
%
%
%
%
%
%
%
%


%
%
%
%

%

\author[chiesa]{Alessandro Chiesa\thanks{Supported by the Center for Long-Term Cybersecurity at UC Berkeley.}}
\author[manohar]{Peter Manohar\thanks{Supported by the NSF Graduate Research Fellowship Program (under Grant No.\ DGE1745016);
and the ARCS Foundation.
Any opinions, findings, and conclusions or recommendations expressed in this material are those of the author(s) and do not necessarily reflect the views of the National Science Foundation.}}
\author[shinkar]{Igor Shinkar\thanks{Supported by NSERC Discovery Grant.}}

%
%
%
%

%
%
%
%
%
%
\begin{abstract}
Many low-degree tests examine the input function via its restrictions to random hyperplanes of a certain dimension. Examples include the line-vs-line (Arora, Sudan 2003), plane-vs-plane (Raz, Safra 1997), and cube-vs-cube (Bhangale, Dinur, Navon 2017) tests.

In this paper we study tests that only consider restrictions along \emph{axis-parallel} hyperplanes, which have been studied by Polishchuk and Spielman (1994) and Ben-Sasson and Sudan (2006).
While such tests are necessarily ``weaker,'' they work for a more general class of codes, namely tensor product codes.
Moreover, axis-parallel tests play a key role in constructing LTCs with inverse polylogarithmic rate and short PCPs (Polishchuk, Spielman 1994; Ben-Sasson, Sudan 2008; Meir 2010). We present two results on axis-parallel tests.

\smallskip
\emph{(1) Bivariate low-degree testing with
%
low agreement.}   %
We prove an analogue of the Bivariate Low-Degree Testing Theorem of Polishchuk and Spielman in the low-agreement regime, albeit
%
for much larger fields.    %
Namely, for the
%
%
tensor product of the Reed--Solomon code with itself,                       %
we prove that for sufficiently large fields, the $2$-query variant of the axis-parallel line test (row-vs-column test) works for \emph{arbitrarily small agreement}. Prior analyses of axis-parallel tests assumed high agreement, and no results for such tests in the low-agreement regime were known.

Our proof technique deviates significantly from that of Polishchuk and Spielman, which relies on algebraic methods such as B\'{e}zout's Theorem, and instead leverages a fundamental result in extremal graph theory by
%
K\H{o}v\'{a}ri,    %
S\'{o}s, and Tur\'{a}n. To our knowledge, this is the first time this result is used in the context of low-degree testing.

\smallskip
\emph{(2) Improved robustness for tensor product codes.}
Robustness is a strengthening of local testability that underlies many applications. We prove that the axis-parallel hyperplane test for the
%
$m$-th tensor power                                                                 %
of a linear code with block length $\blocklength$ and distance $\codedist$ is $\Omega({\codedist^m}/{\blocklength^m})$-robust. This improves on a theorem of Viderman (2012) by a factor of $1/\poly(m)$. While the improvement is not large, we believe that our proof is a notable simplification compared to prior work.
\end{abstract}

\iffalse %
%
%
%
%
%
%
%
%
%
%
%
%
%
%
%
%
%
%
%
%
%

\tocarxivcategory{Computational Complexity}

\fi %

\end{frontmatter}

\section{Introduction}\label{sec:intro}
%


Locally testable codes (LTCs) are error-correcting codes for which, given an input word, one can verify whether the word belongs to or is far from the code by inspecting the word in a few random locations.
LTCs have been studied extensively in different contexts, including program checking, interactive proofs, and probabilistically checkable proofs (PCPs) \cite{BlumLR93,RubinfeldS96,BabaiFL91,BabaiFLS91,PolishchukS94, GoldreichS06}. Goldreich and Sudan~\cite{GoldreichS06} describe LTCs as ``combinatorial counterparts of the complexity theoretic notion of PCPs,'' motivating the study of these objects separately.

\parhead{LTC constructions}
The first constructions of LTCs were algebraic in nature, and relied on multivariate polynomials. Starting with the seminal
%
paper by  %
Blum, Luby, and Rubinfeld~\cite{BlumLR93}, there has been much work on such algebraic LTCs by way of results on \emph{linearity testing} and \emph{low-degree testing} in numerous settings \cite{BlumLR93,BenOrCLR08,BenSassonSVW03,BellareCHKS96,AlonKKLR05,BhattacharyyaKSSZ10}. Many other constructions \cite{Meir09, Viderman13, KoppartyMRS16} further optimize parameters of these codes,
including rate, distance, and the number of queries made by the tester.

Ben-Sasson and Sudan \cite{Ben-SassonS06} suggested a \emph{combinatorial} approach to construct LTCs starting from any linear code by
\begin{inparaenum}[(i)]
 \item applying the \emph{tensor product} operation \cite{Wolf65,WolfE63} to the code, and
 \item testing the resulting code via the \emph{axis-parallel hyperplane test}.
\end{inparaenum}
We now discuss both.

%
%
The tensor product of a linear code $C \subseteq \Field^{\blocklength}$ with itself,    %
denoted $C^{2}$, is the code in $\Field^{\blocklength^{2}}$ consisting of all $2$-dimensional matrices whose $\blocklength$ rows and $\blocklength$ columns are codewords in $C$; similarly,
%
the $m$-th tensor power                                                         %
of $C$, denoted $C^{m}$, is the code in $\Field^{\blocklength^m}$ consisting of all $m$-dimensional
%
tensors     %
$\testerinput$ whose restrictions to any $(m-1)$-dimensional axis-parallel hyperplane is a codeword in $C^{m-1}$. For example, the code of evaluations of all $m$-variate polynomials of individual degree at most $\degree$
%
is the $m$-th tensor power                                                      %
of the code of evaluations of all univariate polynomials of degree at most $\degree$.

The axis-parallel hyperplane test for the code $C^{m}$ works as follows: given a word $\testerinput$, sample a random axis-parallel hyperplane and check if the restriction of $\testerinput$ to this hyperplane is a codeword in $C^{m-1}$. This natural test extends ideas of axis-parallel line tests used in early PCP constructions~\cite{BabaiFL91,BabaiFLS91, AroraS98} to arbitrary tensor product codes.

We study two aspects of the axis-parallel hyperplane test for tensor product codes.

\parhead{(1) Low-agreement regime}
All of the aforementioned
%
papers  %
study the axis-parallel hyperplane test in the ``high-agreement regime,'' in which the given input word is within the unique decoding radius of the tensor product code. What can be said about the ``low-agreement regime,'' in which the given input word may be as far as the list-decoding radius? This setting is more challenging because one wishes to deduce that a given word has some noticeable global correlation with a codeword (or a short list of codewords) by only assuming that a noticeable fraction of local views of the test are local views of (potentially different) codewords.

Results in the low-agreement regime are known for \emph{other} tests, such as tests for the Hadamard code \cite{BellareCHKS96} and the long code \cite{Hastad01} as well as random \emph{non}-axis-parallel hyperplane tests in various dimensions \cite{RazS97,AroraS03,MoshkovitzR08-ldt}. Moreover, these have applications to PCP constructions and hardness of approximation. However, to our knowledge,
%
\emph{prior to our work, no results %
have been  %
known for the low-agreement regime of axis-parallel tests}.

\parhead{(2) Robustness}
Ben-Sasson and Sudan \cite{Ben-SassonS06} analyze the axis-parallel hyperplane test via the notion of \emph{robustness}, a stronger notion of local testability borrowed from the PCP literature \cite{BenSassonGHSV06,DinurR04}. Informally, a test for a code is robust if, given any input that is far from the code, the local view of the test is also far from an accepting view on average. For example, the axis-parallel hyperplane test is robust if, given any $\testerinput$ that is far from $C^{m}$, the restriction of $\testerinput$ to a random hyperplane is far from $C^{m-1}$ on average.

Robustness thus relates the global distance to the expected local view distance and, as shown in \cite{Ben-SassonS06}, facilitates query reduction via a natural way to compose
%
tests.  This notion   %
has also found applications to proof composition in the setting of PCPs \cite{BenSassonGHSV06}.
%
These articles have %
motivated the study of the robustness of the axis-parallel hyperplane test for tensor product codes, establishing both positive results \cite{DinurSW06,BenSassonV08,BenSassonV15,Viderman12} and limitations \cite{Valiant05,CoppersmithR05,GoldreichM12}.

Despite significant progress, robustness results for the axis-parallel hyperplane test seem to be \emph{far from tight}. The best known relation between the global distance and the local distance is due to Viderman \cite{Viderman12}, but no examples that come anywhere close to his proven bound are known.


%
%
%
\section{Main results}
\label{sec:our-results}

We present two main results about tests for tensor product codes. First, we prove an analogue of the Bivariate Low-Degree Testing Theorem of Polishchuk and Spielman \cite{PolishchukS94} in the low-agreement regime, albeit
%
for much larger fields.   %
Second, we improve on the robustness of the hyperplane test for testing the tensor product code $C^{m}$, for $m \geq 3$. We now discuss our results.

%
%
\subsection{Bivariate low-degree testing in the low-agreement regime}
\label{sec:bivariate-testing}

One of the applications of locally testable codes is constructing PCPs, where it is often desirable to reduce the number of queries made by the test. Typically this is done by increasing the alphabet size so that each ``large'' symbol bundles together several ``small'' symbols from different locations of the given word. This bundling now introduces a \emph{consistency} problem, because two large symbols may in principle disagree about the same location in the word.

For example, in \cite{RazS97,AroraS03,MoshkovitzR08-ldt,BhangaleDN17} the test has access to (alleged) restrictions of a low-degree polynomial to all lines, planes, cubes, or other low-degree manifolds. The test samples several queries that intersect, and checks that their answers are consistent on the intersection. These works establish that if the test accepts with probability above a certain threshold, then the restrictions are close to the restrictions of some low-degree polynomial.

We study this problem in a modified setting, where the test only has access to \emph{axis-parallel} restrictions. Restricting the test in this way makes its task more difficult, but doing so provides other advantages. First, axis-parallel restrictions are sometimes the only natural restrictions, such as when testing the
%
$m$-th tensor power                                                                 %
of a general linear code $C$ (more generally, one may consider restrictions to all $(m-1)$-dimensional axis-parallel hyperplanes). Second, having fewer restrictions enables more efficiency, \eg, it facilitates the construction of short PCPs \cite{PolishchukS94,BenSassonS08}.

Indeed, for this very reason, Polishchuk and Spielman \cite{PolishchukS94} study the above problem for bivariate polynomials, where $m=2$ and $C$ is the degree-$\degree$ Reed--Solomon code, \ie, the code of all evaluations $p \colon \Field^2 \to \Field$ of bivariate polynomials of individual degree at most $\degree$. The test has access to a table of row polynomials and a table of column polynomials, and the goal is to check if these tables are consistent with the restrictions of some \emph{global} bivariate polynomial of individual degree $\degree$. That is, for each row $y$, the test gets a polynomial $\rowpoly(\cdot,y) \colon \Field \to \Field$ of degree at most $\degree$ in the first variable $x$. Analogously, for each column $x$, the test gets a polynomial $\colpoly(x,\cdot) \colon \Field \to \Field$ of degree at most $\degree$ in the second variable $y$. The test works as follows: pick a random $(x,y) \in \Field^2$, read the row and column polynomials through this point, and accept if and only if the two polynomials are equal on $(x,y)$.

Clearly, if all the row polynomials and column polynomials are restrictions of a bivariate polynomial of individual degree $\degree$, then the test always accepts. They prove that, conversely, if the test accepts with probability close to 1, then the given polynomials are ``close'' to being restrictions (to axis-parallel lines) of some low-degree bivariate polynomial, as written below. In the statement, we say that a bivariate polynomial in variables $x$ and $y$ has degree $(a,b)$ if the degree in $x$ is at most $a$ and that in $y$ is at most $b$. This means that the table of row polynomials, $\rowpoly(x,y)$, has degree $(\degree, \blocklength)$ and the table of column polynomials, $\colpoly(x,y)$, has degree $(\blocklength, \degree)$, where $\blocklength$ is the size of the table.

\begin{theorem}[\cite{PolishchukS94}]
\label{thm:PS94}
Let $\Field$ be a field and $X,Y \subseteq \Field$ subsets of size $\blocklength \DefineEqual \abs{X} = \abs{Y}$.
Let $\rowpoly(x,y)$ be a polynomial of degree $(\degree,\blocklength)$ and $\colpoly(x,y)$ a polynomial of degree $(\blocklength,\degree)$ such that
\begin{equation*}
\Pr_{(x,y) \in X \times Y}
[\colpoly(x,y) = \rowpoly(x,y)]
= 1 - \PSprox^2
\end{equation*}
for some $\PSprox>0$.
If $\blocklength > 2 \PSprox \blocklength + 2\degree$, then there exists a polynomial $Q(x,y)$ of degree $(\degree,\degree)$ such that
\begin{equation*}
\Pr_{(x,y) \in X \times Y}[\colpoly(x,y) = \rowpoly(x,y) = Q(x,y)] \geq 1 - 2\PSprox^2 \enspace.
\end{equation*}
\end{theorem}

The theorem above assumes that $\blocklength > 2 \PSprox \blocklength + 2\degree$, which means that $\PSprox^2 < (1/2-\degree/\blocklength)^{2} < 1/4$. In other words, it requires the row polynomials and column polynomials to agree on (at least) more than three quarters of the points in $X \times Y$.
A slight improvement in the parameters of this theorem is shown in \cite{BenSassonCGT13}. However, their result still requires the polynomials to agree on a large fraction of the points in $X \times Y$.
But what, if anything, can be said if we only assume that they agree, for example, on more than a $0.1$-fraction of those points?

There are several results on low-degree testing that show that, even if we only assume that the test accepts with noticeable probability (for the row-vs-column test this probability equals the agreement between row and column polynomials), one can \emph{still} prove the existence of a \emph{short list} of polynomials that ``explain'' most of this probability, and this in turn has applications to constructing PCPs with small errors (see, \eg,~\cite{RazS97,AroraS03,MoshkovitzR08-ldt}).

Our next result gives a positive answer to the question above, stating that even in the low-agreement regime, we can still deduce some structure about the polynomials $\rowpoly$ and $\colpoly$, assuming that the
%
order of the field  %
is sufficiently large.

\begin{ftheorem}
\label{thm:bivar-low-agreement}
Fix $\degree \in \Naturals$, and let $\prox, \eps \in \mathbb R$ be such that $\prox \geq \eps > 0$. Let $\Field$ be a field of
%
order  %
$\blocklength$, where $\blocklength > \exp(\Omega(({\degree}/{\eps}) \log({1}/{\eps})))$. Let $\rowpoly(x,y)$ be a polynomial of degree $(\degree,\blocklength)$ and $\colpoly(x,y)$ a polynomial of degree $(\blocklength,\degree)$ such that
\begin{equation*}
\Pr_{(x,y) \in \Field^2}[\colpoly(x,y) = \rowpoly(x,y)] = \prox \enspace.
\end{equation*}
Then there exist $t = O({1}/{\eps})$ polynomials $Q_1(x,y),\dots,Q_t(x,y)$ of degree $(\degree,\degree)$ such that
\begin{equation*}
\Pr_{(x,y) \in \Field^2}
[\exists i \in [t] \ \colpoly(x,y) = \rowpoly(x,y)= Q_i(x,y)]
\geq \prox - \eps \enspace.
\end{equation*}

\end{ftheorem}

We remark that \expref{Theorem}{thm:bivar-low-agreement} holds in general
%
%
for $C^2$, where $C \subseteq \Field^{\blocklength}$ is \emph{any} linear code      %
with minimal distance $\geq \blocklength-\degree$ such that
$\blocklength > \exp(\Omega(({\degree}/{\eps}) \log({1}/{\eps})))$.
In particular, this means that the minimal distance of $C$ is at least $\blocklength - O(\log \blocklength)$.
See the paragraph \emph{Beyond polynomials} on page~\pageref{par:beyondpoly} for details.

Note that in the above theorem, $\prox$ is the agreement probability, while $\PSprox^{2}$ in \expref{Theorem}{thm:PS94} is the disagreement probability. Also, both $\prox$ and $\eps$ can be \emph{sub-constant}. This is the first result that analyzes the row-vs-column test
%
in the low-agreement regime that we are aware of.              %

The row-vs-column test and its higher-dimensional analogues underlie many known PCP constructions \cite{BabaiFL91, BabaiFLS91, PolishchukS94, BenSassonS08}.
However, in all these constructions the
%
low-degree tests are only analyzed in the
%
high-agreement regime.   %
We believe that analyzing the test in the low-agreement regime may imply short PCP constructions with small (sub-constant) soundness.
A weakness of the result stated in \expref{Theorem}{thm:bivar-low-agreement} is the requirement that the field
%
must be very large,
which restricts us from getting PCPs with
%
polynomial-length proofs.  %
Nonetheless, we consider \expref{Theorem}{thm:bivar-low-agreement} as a promising first step in this direction.
More generally, our result suggests that the low-agreement regime for tensor product codes merits further study.


To prove the theorem we leverage a fundamental result in extremal graph theory by %
K\H{o}v\'{a}ri,    %
S\'{o}s, and Tur\'{a}n.
To our knowledge, this is the first time this result is used in the context of low-degree testing.
See \expref{Subsection}{sec:techniques-2} below for a high-level description of our proof.

We observe that in the above theorem, for $\prox \geq \eps$ sufficiently small, it is \emph{necessary} to have a
list of at least ${(\prox/\eps)}/{\polylog(\prox/\eps)}$ polynomials in order to explain all but $\eps$ of the agreements of $\rowpoly$ and $\colpoly$.
In particular, if $\prox$ is a small constant and $\eps>0$ is sufficiently small,
then the bound in \expref{Theorem}{thm:bivar-low-agreement} is tight up to a $\polylog(1/\eps)$ factor.

\begin{proposition}
\label{prop:list size lb}
Fix $\degree \in \Naturals$, and let $\prox, \eps \in \mathbb R$ be such that ${1}/{12} > \prox \geq \eps > 0$. Let $\Field$ be a field of
%
order   %
$\blocklength$, where
\[
  \blocklength > \Omega\left( \frac{\degree}{\eps} \cdot \frac{\delta}{\eps} \polylog({1}/{\eps})\right)\enspace.
\]
Then there exists a polynomial $\rowpoly(x,y)$ of degree
$(0,\blocklength)$ and a polynomial $\colpoly(x,y)$ of degree
$(\blocklength,0)$ such that
\begin{equation*}
\Pr_{(x,y) \in \Field^2}[\colpoly(x,y) = \rowpoly(x,y)] \in [\prox,2\prox] \enspace,
\end{equation*}
but for any $t < {(\prox/\eps)}/{\polylog(\prox/\eps)}$ and polynomials $Q_1(x,y),\dots,Q_t(x,y)$ of degree $(\degree,\degree)$ it holds that
\begin{equation*}
\Pr_{(x,y) \in \Field^2}
[\exists i \in [t] \st \colpoly(x,y) = \rowpoly(x,y)= Q_i(x,y)]
< \Pr_{(x,y) \in \Field^2}[\colpoly(x,y) = \rowpoly(x,y)] - \eps \enspace.
\end{equation*}
\end{proposition}

Thus, unlike other results in this area \cite{RazS97, AroraS03, BhangaleDN17} where the list size depends only on $\prox$, the list size \emph{must} grow as a function of $\eps$ in our setting. The reason for this difference comes from the restriction of our test, which considers only axis-parallel lines, as opposed to arbitrary lines (or planes or cubes) used in the other works. In particular, in our setting once we choose a point $(x,y)$, the lines going through this point are fixed by the design of the test, while in the
%
aforementioned papers  %
there are many lines (or planes or cubes) going through this point, and the performed queries are chosen at random conditioned on the chosen random point.


%
%
%
%
%
%
%
%
\parhead{Beyond polynomials}\label{par:beyondpoly}
While
%
the proof in~\cite{PolishchukS94}   %
relies on polynomials (a key step is B\'{e}zout's Theorem), we rely on combinatorial techniques, so that our \expref{Theorem}{thm:bivar-low-agreement} holds in general
%
%
for $C^2$, where $C \subseteq \Field^{\blocklength}$ is \emph{any} linear code      %
with minimal distance $\geq \blocklength-\degree$ such that $\blocklength > \exp(\Omega(({\degree}/{\eps}) \log({1}/{\eps})))$.
In particular, this means that the minimal distance of $C$ is at least $\blocklength - O(\log \blocklength)$.
The row-vs-column test is now given two matrices $\rowpoly,\colpoly \in \Field^{\blocklength \times \blocklength}$ such that every row of $\rowpoly$ is in $C$, and every column of $\colpoly$ is in $C$. If
\[
  \Pr_{(x,y) \in [\blocklength]^2}[\rowpoly(x,y) = \colpoly(x,y)] = \prox\enspace,
\]
then there exist $t = O(1/\eps)$ codewords $Q_1,\dots,Q_t \in
C^2$, %
such that
\[
  \Pr_{(x,y) \in [\blocklength]^2}[\exists i \in [t] \st \rowpoly(x,y)
  = \colpoly(x,y) = Q_i(x,y)] \geq \prox - \eps\enspace.
\]


In this context it is worth mentioning that there has been a lot of work on the robustness of the axis-parallel line test for
%
tensor products of linear codes,                                                %
proving both positive results \cite{Ben-SassonS06,DinurSW06} and negative ones \cite{Valiant05,CoppersmithR05,GoldreichM12}.
We find it quite remarkable that this result holds for
%
the tensor product $C^2$ of an arbitrary code $C$,                                          %
albeit with very high distance,
as the closely related notion of robustness does not hold for
%
the tensor product $C^2$ of an arbitrary code $C$.                                          %
Finally, in the high-agreement regime there is a \emph{correspondence} between
the robustness of the axis-parallel line test
and the soundness of the row-vs-column test (the matrix is given as a collection of lines rather than explicitly).%
\footnote{
Let $M \in \Field^{\blocklength \times \blocklength}$ be such that
the average relative distance of a row/column of $M$ to some codeword is $1-\eps$.
One can verify that by considering the closest codewords in each row and in each column,
the obtained table of row/column codewords
passes the row-vs-column test with probability at least $1-2\eps$.
Therefore, there exists a tensor codeword that agrees with most of the rows and most of the columns,
which in turn implies its agreement with $M$.}
Yet this correspondence \emph{does not hold} in the low-agreement regime. Consider a matrix $M$ whose rows are codewords chosen uniformly at random independently from each other: assuming that $\Field$ is sufficiently large, the average relative distance of a row/column of $M$ from some codeword is approximately ${1}/{2} + {k}/{(2\blocklength)}$, where $k$ is the dimension of the code. This is because when reading a row the test reads a codeword of $C$, and when reading a column, then the view agrees with a codeword on approximately $k$ coordinates with high probability.
On the other hand, $M$ is typically far from a tensor codeword, \eg, it agrees with any tensor codeword on at most $k/n$ fraction of points.

\parhead{Open problems}

We raise two questions on the low-agreement regime of axis-parallel line tests.
\begin{itemize}

  \item \emph{Smaller block length.}
  Our result (\expref{Theorem}{thm:bivar-low-agreement}) assumes that the length of code $\blocklength$ is exponential in the degree $\degree$. Can one prove a similar result for smaller block length, \eg, $\blocklength=\poly(\degree)$?

  \item \emph{Higher dimensions.}
  Polishchuk and Spielman \cite{PolishchukS94} explain that their result (in the high-acceptance regime) also holds in higher dimensions, where now the test is given a table of low-degree polynomials for each axis-parallel line in $\Field^m$ and works as follows: pick a random $p \in \Field^m$, read the polynomials along the $m$ axis-parallel lines through $p$, and check that all polynomials agree on $p$. Can one prove a high-dimensional analogue of \expref{Theorem}{thm:bivar-low-agreement}? Namely, is it true that if this test accepts with probability $\prox > 0$, then there is a short list of low-degree polynomials that explain most of the agreements?

\end{itemize}
%
%
\subsection{Improved robustness for the axis-parallel hyperplane test}
\label{sec:result-robustness}

We study the robustness of the axis-parallel hyperplane test for the tensor product code $C^{m} \subseteq \Field^{\blocklength^{m}}$, for an arbitrary linear code $C$ with minimal distance $\codedist$ and block length $\blocklength$ over the field $\Field$. Let $\tester$ be the test that, given a word $\testerinput \in \Field^{\blocklength^{m}}$, samples a random $(m-1)$-dimensional axis-parallel hyperplane $\plane$ and checks if $\testerinput \vert_{\plane} \in C^{m-1}$. For a word $\testerinput \in \Field^{\blocklength^{m}}$, we define $\reldist(\testerinput)$ to be the relative distance of the word $\testerinput$ to the code $C^{m}$ and $\robust(\testerinput)$ to be $\mathbb E_{\plane} [\reldist(\testerinput \vert_{\plane}, C^{m-1})]$, the expected local distance of $\testerinput$. The test $\tester$ is \emph{$\alpha$-robust} if $\robust(\testerinput) \geq \alpha \cdot \reldist(\testerinput)$ for every word $\testerinput \in \Field^{\blocklength^{m}}$. The ``strength'' of the test increases with $\alpha$, so the goal is to establish the largest $\alpha$ for which this inequality holds.

\parhead{What is known}
There are two main
%
pieces of prior work  %
that study the robustness of the test $\tester$ for general $m$.
We state the results of these
%
articles  %
starting with one of Ben-Sasson and Sudan \cite{Ben-SassonS06}.

\begin{theorem}[\cite{Ben-SassonS06}]
\label{thm:bs}
Let $C \subseteq \Field^n$ be a linear code with minimal distance $d$.
For
\[
  m \geq 3\qquad\text{and}\qquad \left(\frac{\codedist-1}{\blocklength}\right)^m \geq 7/8\enspace,
\]
the test $\tester$ is $\alpha$-robust for $C^m$ with $\alpha = 2^{-16}$.
\end{theorem}

\noindent
The above theorem is limited in that the proved robustness is small and, moreover, only provides a guarantee when $C$ has a very large distance. Viderman \cite{Viderman12} shows that this condition on the distance is not necessary in order to show \emph{some} robustness guarantee.

\begin{theorem}[\cite{Viderman12}]
\label{thm:vid}
Let $C \subseteq \Field^n$ be a linear code with minimal distance $d$.
For $m \geq 3$, the test $\tester$ is $\alpha$-robust %
for   %
$C^m$ with
\[
  \alpha = \frac{1}{2 m^2} \left(\frac{\codedist}{\blocklength}\right)^m\enspace.
\]
\end{theorem}

\noindent
The above theorem, the state of the art in this setting, improves on the previous one as
\begin{inparaenum}[(i)]
\item even if
  \[
    \left(\frac{\codedist-1}{\blocklength}\right)^{m} \geq 7/8\enspace,
  \]
  the robustness provided by \expref{Theorem}{thm:vid} is larger than
  that provided by \expref{Theorem}{thm:bs} for $m \leq 169$;
  \item a robustness guarantee is provided for any choice of $m,\codedist,\blocklength$ (as long as $m \geq 3$).
\end{inparaenum}

\parhead{Our result}
We present a simpler proof of \expref{Theorem}{thm:vid}, which also achieves a $m^2$ improvement in the robustness by showing that the hyperplane test is $\Omega({\codedist^m}/{\blocklength^m})$-robust. This improved value for the robustness appears more ``natural,'' because ${\codedist^m}/{\blocklength^m}$ is the distance of the code $C^m$.

\begin{ftheorem}
\label{thm:pat-robustness}
Let $C \subseteq \Field^n$ be a linear code with minimal distance $d$.
For $m \geq 3$, the test $\tester$ is $\alpha$-robust for $C^m$ with
\[
  \alpha = \frac{1}{12} \left(\frac{\codedist}{\blocklength}\right)^m\enspace.
\]
\end{ftheorem}

%
%


\parhead{Tight or not?}
%
%
The test $\tester$ has been studied in several papers,          %
and all resulting analyses have an exponential dependence on $m$ in the robustness. Yet, there is no evidence indicating that this dependence is necessary. Perhaps a ``dream'' result of constant robustness, for all codes $C$ and $m \geq 3$, is possible. Like previous results, we too incur the same exponential dependence in the robustness. We present some observations that may suggest that this dependence is not necessary.


%
\begin{itemize}

\item Under certain conditions on $\testerinput$, we can prove that \[
    \robust(\testerinput) \geq \max\left\{\frac{1}{m + c},\ c' \frac{\codedist^m}{\blocklength^m}\right\} \cdot \reldist(\testerinput)
  \]
  for constants $c,c'>0$. These two expressions are
  \emph{incomparable}, as we can set the parameters
  $m,\codedist,\blocklength$ to make either expression bigger than the
  other. (See \expref{Claim}{claim:m}.)

\item The guarantees of \expref{Theorem}{thm:bs},
  \expref{Theorem}{thm:vid}, and \expref{Theorem}{thm:pat-robustness}
  all degrade as ${\codedist^m}/{\blocklength^m}$ decreases. In
  particular, the proven value of $\alpha$ in all these cases tends to
  $0$ as ${\codedist}/{\blocklength}$ tends to $0$. However, for
  any code $C$ we can prove that
  \[
    \reldist(\testerinput) \leq \robust(\testerinput) + \frac{\blocklength - k}{\blocklength}\enspace.
  \]
  In particular, we show that if $C$ is an MDS code, then
  \[
    \reldist(\testerinput) \leq \robust(\testerinput) + \frac{\codedist
      - 1}{\blocklength}
  \]
  for all $\testerinput$. (See
  \expref{Corollary}{cor:robustness-of-RS}.)

\end{itemize}
We, thus, think that determining the optimal robustness of $\tester$ is an intriguing open problem:
\begin{center}
\begin{minipage}{0.9\textwidth}
  What is the optimal robustness of the hyperplane test $\tester$? \\ Can one prove that
  \[
    \alpha = \Omega\left(\max\left(\frac{1}{m},\ \frac{\codedist^m}{\blocklength^m}\right)\right)\enspace,
  \]
  or even $\alpha = \Omega(1)$, for all codes?
\end{minipage}
\end{center}
In \cite{Ben-SassonS06}, \cite{Viderman12}, and our result, the proof shows that when $\robust(\testerinput)$ is below some threshold (related to the code's unique decoding radius), then $\reldist(\testerinput)$ is also small. However, when $\robust(\testerinput)$ is not below this threshold, the analysis says nothing about $\reldist(\testerinput)$, and naively uses $\reldist(\testerinput) \leq 1$ to prove robustness in this regime. We believe that progress on understanding the optimal robustness of $\tester$ hinges on understanding what techniques (if any) can be used to bound $\reldist(\testerinput)$ in terms of $\robust(\testerinput)$ for a larger range of $\robust(\testerinput)$.

%
%
%
\subsection{Roadmap}
\label{sec:roadmap}

The rest of this paper is organized as follows.
\expref{Section}{sec:techniques} describes the techniques used to prove our results.
\expref{Section}{sec:prelim-bivar} introduces preliminaries for \expref{Theorem}{thm:bivar-low-agreement}.
\expref{Section}{sec:proof2} provides the proof of \expref{Theorem}{thm:bivar-low-agreement} and  \expref{Proposition}{prop:list size lb} (which justifies the list size in \expref{Theorem}{thm:bivar-low-agreement}).
\expref{Section}{sec:prelim-pat} introduces basic notations and definitions for \expref{Theorem}{thm:pat-robustness}.
\expref{Section}{sec:proof1} provides the proof of \expref{Theorem}{thm:pat-robustness}.



%
%
%
\section{Techniques}
\label{sec:techniques}

We give an overview of the proof techniques behind  \expref{Theorem}{thm:bivar-low-agreement} and \expref{Theorem}{thm:pat-robustness}.

%
%
%
\subsection{\expref{Theorem}{thm:bivar-low-agreement}: bivariate testing in the low-agreement regime}       %
\label{sec:techniques-2}

Polishchuk and Spielman~\cite{PolishchukS94} prove their result (\expref{Theorem}{thm:PS94}) using the following approach. Given $\rowpoly$ and $\colpoly$ (as in the theorem) such that $\Pr_{x,y}[\rowpoly(x,y) = \colpoly(x,y)] > 1 - \prox$, they define an ``error polynomial'' $E$ that equals 0 for all $(x,y)$ such that $\rowpoly(x,y) \neq \colpoly(x,y)$. Since the fraction of points where $\rowpoly(x,y) \neq \colpoly(x,y)$ is small, $E$ is a low-degree polynomial. However, in the low-agreement regime that we consider, the degree of $E$ is rather large, which seems to preclude their approach. In particular, a key step based on B\'{e}zout's Theorem in their proof appears to break down.

We take a completely different approach, which relies on a combinatorial statement from extremal graph theory. Given $\rowpoly$ and $\colpoly$ such that $\Pr_{x,y}[\rowpoly(x,y) = \colpoly(x,y)] = \prox$, we define $A \in \{0,1\}^{\blocklength \times \blocklength}$ to be the ``agreement matrix'': $A(x,y) = 1$ if and only if $\rowpoly(x,y) = \colpoly(x,y)$. By the assumption it follows that $A$ has at least $\prox \blocklength^2$ ones. By invoking the
%
K\H{o}v\'{a}ri--S\'{o}s--and Tur\'{a}n
Theorem (which may be thought of as an analogue of Tur\'{a}n's Theorem for bipartite graphs) it follows that there are some $S,T \subseteq [\blocklength]$ such that $\abs{S} , \abs{T} > \Omega(\log(\blocklength)) \gg \degree$ and $A\vert_{S \times T} \equiv 1$. Since the rows of $\rowpoly$ and the columns of $\colpoly$ are polynomials of degree $\degree$, we deduce that there exists a unique polynomial $Q$ of degree $(\degree,\degree)$ such that for all $(x,y) \in S \times T$ it holds that $\rowpoly(x,y) = \colpoly(x,y) = Q(x,y)$.

The argument above may appear to be good progress toward our goal. However, there is a total of $\approx \prox \blocklength^2$ ones in $A$, and the rectangle $S \times T$ is of size $O(\log(\blocklength))$, \ie, tiny compared to $\blocklength$. This means that the progress is actually rather small!

Nevertheless, we can now set $A\vert_{S \times T}$ to be zero, and repeat the same argument again, thus covering all but a small fraction of ones of $A$ with small rectangles. However, this raises a new problem. Each rectangle $S \times T$ found in the previous step can be \emph{very} small, and so there are potentially many different polynomials $Q$ that explain the agreements of $\rowpoly$ and $\colpoly$. Our next goal is therefore to ``stitch'' these rectangles together to show that, in fact, there is only a \emph{small} number of distinct polynomials. We do so by ``making the rectangles larger,'' as we now explain.

Consider a rectangle $S \times T$ from the first step, and let $t' \in \Field \setminus T$. Note that if there are $\degree+1$ points $s' \in S$ such that $A(s',t') =1$, then the row polynomial $\rowpoly( \cdot, t')$ is uniquely defined by these $\degree+1$ points, and hence $A(s,t') = 1$ for all $s \in S$. Therefore, we can increase $T$ by adding $t'$ to it. On the other hand, if there are less that $\degree+1$ such points $s'\in S$, then we may disregard these points as they amount to only a small fraction of the points (since $\abs{S} \gg \degree$). Thus, on a typical rectangle $S \times T$, we can go from size $O(\log(\blocklength)) \times O(\log(\blocklength)) $ to size roughly $O(\log(\blocklength)) \times \Omega(\blocklength)$.

In the last step, we show that if we have many rectangles of size $O(\log(\blocklength)) \times \Omega(\blocklength)$ then it is possible to ``stitch'' them together using the fact that if we have two rectangles $S_1 \times T_1$ and $S_2 \times T_2$ with corresponding polynomials $Q_1$ and $Q_2$ such that $\abs{T_1 \cap T_2} > \degree$, then $Q_1 \equiv Q_2$. This follows by the fact that if two univariate polynomials of degree $\degree$ agree on more than $\degree$ points, then they are equal. We then use the inclusion-exclusion principle and some simple inequalities to show that we cannot have more than ${2}/{\eps}$ subsets $T_i \subseteq [n]$ of size at least $\eps\blocklength$ such that $\abs{T_i \cap T_j} \leq \degree$ for all $i \neq j$.

The full proof of \expref{Theorem}{thm:bivar-low-agreement} is provided in \expref{Section}{sec:proof2}.

%
%
\subsection{\expref{Theorem}{thm:pat-robustness}: improved robustness for the hyperplane test}
\label{sec:techniques-1}

Our goal is to prove that the axis-parallel hyperplane test $\tester$ is $\alpha$-robust for
\[
  \alpha = \frac{1}{12} \left(\frac{\codedist}{\blocklength}\right)^m\enspace.
\]
We prove this statement via a careful combination of the approaches
taken by \cite{Ben-SassonS06} and \cite{Viderman12}. Specifically, we
analyze $\robust(\testerinput)$ and $\reldist(\testerinput)$ by
studying the following combinatorial object: the \emph{inconsistency
  graph $G$} of the hyperplane test $\tester$, which we now informally
describe.

The test $\tester$ has access to a word $\testerinput \in \Field^{\blocklength^m}$, allegedly in $C^{m}$. For any axis-parallel hyperplane $\plane$, we denote by $\closestlocalword{\plane}$ the closest codeword to $\testerinput \vert_{\plane}$ in $C^{m-1}$ (breaking ties by picking an arbitrary closest codeword). The vertex set of the graph $G$ is the set of $(m-1)$-dimensional axis-parallel hyperplanes, which are the local views of the test. There is an edge between two different hyperplanes $\plane$ and $\otherplane$ if $\closestlocalword{\plane}$ and $\closestlocalword{\otherplane}$ disagree on the intersection of the hyperplanes, $\plane \cap \otherplane$. (See \expref{Definition}{def:graph} for details.) In other words, the graph has an edge between two planes if the local codewords assigned to the planes are inconsistent. The graph $G$ that we study is similar to the inconsistency graph analyzed in \cite{Ben-SassonS06}. The difference is that, for some threshold parameter $\tau$, the graph used in \cite{Ben-SassonS06} adds an edge from $\plane$ to every other $\otherplane$ in the graph if $\reldist(\testerinput \vert_{\plane},\closestlocalword{\plane}) > \tau$.

First, we show that if $G$ has a large independent set $I$, then there is a codeword $\closestglobalword$ in $C^{m}$ that agrees with the local codewords $\closestlocalword{\plane}$ on \emph{every} hyperplane $\plane$ in $I$. For an independent set $I$, we define $I_b$ to be the set of $i \in [\blocklength]$ such that the hyperplane $\{p \in [\blocklength]^m : p_b = i\}$ is in $I$. A key property of tensor product codes is the unique extension property, which we formally state later on as \expref{Claim}{claim:tensorext}. Using the unique extension property of tensor product codes, we show that if there are two axes $b_1$ and $b_2$ where $I_{b_1}$ and $I_{b_2}$ both have at least $\blocklength - \codedist + 1$ planes, then there is a word $\closestglobalword$ in $C^m$ where $\closestglobalword \vert_{\plane} = \closestlocalword{\plane}$ for every $\plane$ in the independent set. Without loss of generality assume $b_1 = 1$ and $b_2 = 2$. Intuitively, we fill in the restricted hypercube in $\Field^{I_1 \times I_2 \times \blocklength^{m-2}}$ with the values of the closest codewords to $\testerinput \vert_{\plane}$ for each $\plane$ in the independent set. Since the independent set is large, the restricted hypercube is large enough so that we can extend the partially filled-in hypercube to a unique codeword $\closestglobalword$ in $C^m$. The uniqueness of the extension implies that $\closestglobalword \vert_{\plane} = \closestlocalword{\plane}$ for every $\plane$ in $I$.

Next, we analyze the structure of $G$ to show that every edge is adjacent to a vertex of degree at least $(m-2)\codedist/2$. The key point is that two different $C^{m-2}$ codewords must disagree on at least $\codedist^{m-2}$ points, and these points have a particular structure. For two distinct $C^{m-2}$ codewords, we prove that on each of the $m-2$ remaining axes there must be at least $\codedist$ planes, parallel to that axis, that contain points of disagreement. If not, then using the unique extension property we show that the two codewords must be equal, which is a contradiction. For any edge $(\plane, \otherplane)$, this gives us a total of $(m-2)\codedist$ planes that disagree with at least one of $\closestlocalword{\plane}$ and $\closestlocalword{\otherplane}$ on $\plane \cap \otherplane$, which shows that $\deg(\plane) + \deg(\otherplane)$ is at least $(m-2)\codedist$. Therefore, at least one of $\plane$ and $\otherplane$ has degree at least $(m-2)\codedist/2$. As an immediate consequence,
%
the set of planes of degree                                                         %
at least $(m-2)\codedist/2$, which we denote by $L$, is a \emph{vertex cover}, and the set of planes not in $L$ is an \emph{independent set} $I$.

With some algebraic manipulation, we relate the size of this vertex cover to the expected local distance $\robust(\testerinput)$. By expressing $\robust(\testerinput)$ as a sum over pairs of intersecting planes, we show that
\begin{equation*}
\robust(\testerinput) \geq  \frac{1}{\blocklength^{m} m(m-1)} \sum_{(\plane, \otherplane) : \plane \cap \otherplane \ne \emptyset} \hamdist \vert_{\plane \cap \otherplane} (\closestlocalword{\plane}, \closestlocalword{\otherplane}) \enspace.
\end{equation*}
This allows us to express the robustness of the test $\tester$ in terms of the size of the vertex cover $L$.

Similar to the analysis of \cite{Viderman12}, we break up the proof into two cases. If $\abs{L}$ is somewhat large,
then
\[
  \rho(M) \geq \frac{1}{12} \left(\frac{\codedist}{\blocklength}\right)^m\enspace,
\]
and the theorem follows immediately because $\reldist(\testerinput)$
is anyways at most $1$. If $\abs{L}$ is small, then the corresponding
independent set has two axes where
$\abs{I_b} \geq \blocklength - \codedist + 1$. Therefore, there is a
global codeword $\closestglobalword$ that is consistent with all the
hyperplanes in the independent set. We use this fact to show that
$\reldist(\testerinput)$ must be small when $\robust(\testerinput)$ is
small, which concludes the proof.

The full proof of \expref{Theorem}{thm:pat-robustness} is provided in \expref{Section}{sec:proof1}.

%
%
\section{Preliminaries for \expref{Theorem}{thm:bivar-low-agreement}}
\label{sec:prelim-bivar}
\subsection{Low-degree polynomials}
We will use the following lemmas about low-degree polynomials in the proof of \expref{Theorem}{thm:bivar-low-agreement}. These are standard interpolation lemmas, and direct proofs can be found in \cite{PolishchukS94}.

\begin{lemma}\label{lemma:zero poly}
    Let $S,T \subseteq \Field$ be two sets each of size at least $\degree+1$.
    Suppose that for two polynomials $Q_1(x,y), Q_2(x,y)$ of degree $(\degree,\degree)$,
    if holds that $Q_1(x,y) = Q_2(x,y)$ for all $(x,y) \in S \times T$.
    Then $Q_1 \equiv Q_2$.
\end{lemma}

\begin{lemma}\label{lemma:interpolation}
    Let $S,T \subseteq \Field$ be two sets each of size at least $\degree+1$.
    Suppose that there is polynomial $\rowpoly(x,y)$ of degree $(\degree, \blocklength)$,
    and a polynomial $\colpoly(x,y)$ of degree $(\blocklength,\degree)$ such that
    $\rowpoly(x,y) = \colpoly(x,y)$ for all $(x,y) \in S \times T$.
    Then, there exists a polynomial $Q(x,y)$ of degree $(\degree,\degree)$
    such that $Q(x,y) = \colpoly(x,y) = \rowpoly(x,y)$ for all $(x,y) \in S \times T$.
\end{lemma}

\begin{corollary}\label{cor:completing boxes}
    Let $S,T \subseteq \Field$ be two sets each of sizes $\abs{S} \geq \degree+2$ and $\abs{T} \geq \degree+2$,
    and let $(x_0,y_0) \in S \times T$.
    Suppose that there is a polynomial $\rowpoly(x,y)$ of degree $(\degree,\blocklength)$,
    and a polynomial $\colpoly(x,y)$ of degree $(\blocklength,\degree)$ such that
    $\rowpoly(x,y) = \colpoly(x,y)$ for all $(x,y) \in S \times T \setminus \{(x_0,y_0)\}$.
    Then $\colpoly(x_0,y_0) = \rowpoly(x_0,y_0)$.
\end{corollary}

%
\subsection{The \texorpdfstring{K\H{o}v\'{a}ri}{Kovari}--S\'{o}s--Tur\'{a}n theorem} %

We first define the density of a binary matrix.
\begin{definition}
    Let $A \in \{0,1\}^{k \times \ell}$ be a binary matrix.
    Define the \emph{density} of $A$ to be
    \[
      \delta(A) = \frac{\sum_{i \in [k], j \in [\ell]} A_{i,j}}{k \cdot \ell}\enspace.
    \]
    We say that $A$ is \emph{$\tau$-dense} if $\delta(A) \geq \tau$.
\end{definition}

In the proof of \expref{Theorem}{thm:bivar-low-agreement}
we will use a result due to
%
K\H{o}v\'{a}ri,     %
S\'{o}s, and Tur\'{a}n~\cite{KovariST54},
which states that any sufficiently dense binary matrix contains a large
submatrix where every entry is $1$.

%
\begin{theorem}[K\H{o}v\'{a}ri, S\'{o}s, Tur\'{a}n]\label{thm:KST}  %
Let $N, M, t, s$ be natural numbers that satisfy $N \ge s$ and $M \ge t \ge s$,
and let $A \in \{0,1\}^{N \times M}$ be a binary matrix.
If $A$ is $(\sqrt[s]{{t-1}/{M}} + {s}/{N})$-dense, then
there are $S \subseteq [N]$ and $T \subseteq [M]$ of sizes $\abs{S} = s$ and $\abs{T}=t$
such that $A \vert_{S \times T} \equiv 1$.
\end{theorem}

\begin{remark}
  The %
K\H{o}v\'{a}ri--S\'{o}s--Tur\'{a}n   %
  theorem is usually stated as saying that
    any sufficiently dense bipartite graph contains a large bipartite clique.
    It is clear, however, that the matrix formulation above is equivalent
    by associating a bipartite graph with its adjacency matrix, where the rows correspond to the vertices
    on the left, and the columns correspond to the vertices on the right.
\end{remark}
%
%




%
%
%
\section{Proof of \expref{Theorem}{thm:bivar-low-agreement}}
\label{sec:proof2}

The key step in the proof of \expref{Theorem}{thm:bivar-low-agreement} is the following lemma.

\begin{lemma}[Key lemma]\label{lemma:key}
  Fix $\eps > 0$. Let $\Field$ be a field of
%
  order %
  $\blocklength$, and suppose that
  \[
    \blocklength > \exp\left(\Omega\left(\frac{\degree}{\eps}
    \log({1}/{\eps})\right)\right)\enspace.
  \]
  Then there are $t \leq {2}/{\eps}$ polynomials $Q_1,\dots,Q_t$
  each of degree $(\degree,\degree)$, and subsets
  $S_1,\dots,S_t,B_1,\dots,B_t \subseteq \Field$ such that:
    \begin{enumerate}
        \item
            For all $i \in [t]$ and $(x,y) \in S_i \times B_i$ it holds that $\colpoly(x,y) = \rowpoly(x,y) = Q_i(x,y)$.
        \item
%
%
          The sets $S_i$ are pairwise disjoint.                 %
%
%
        \item
 %
%
  $\displaystyle{\frac{\abs{\bigcup_{i \in [t]} S_i \times B_i}}{\abs{\Field}^2} \geq \prox - 3\eps}$,          %
            where $\prox = \Pr[\colpoly(x,y) = \rowpoly(x,y)]$.
    \end{enumerate}
\end{lemma}

Before proving \expref{Lemma}{lemma:key}, let us see how it immediately implies \expref{Theorem}{thm:bivar-low-agreement}.

\begin{proof}[Proof of \expref{Theorem}{thm:bivar-low-agreement} using \expref{Lemma}{lemma:key}]
    Let $\eps  > 0$, and apply \expref{Lemma}{lemma:key}
    with $\eps/3$.
    By \expref{Lemma}{lemma:key} for some
    $t \leq {2}/{(\eps/3)} = {6}/{\eps}$
    there are disjoint subsets $S_1 \times B_1,\dots,S_t \times B_t \subseteq \Field^2$
    such that
%
%
    \[{\frac{\abs{\bigcup_{i \in [t]} (S_i \times B_i)}}{\abs{\Field}^2} \geq \prox - \eps}\enspace,
    \]
    and %
    for all $i \in [t]$ and $(x,y) \in S_i \times B_i$ it holds that
    $\rowpoly(x,y) = \colpoly(x,y) = Q_i(x,y)$.  This implies that
    \begin{equation*}
        \Pr_{(x,y) \in \Field^2}[\exists i \in [t] \st \rowpoly(x,y) = \colpoly(x,y)= Q_i(x,y)] \geq
%
  \Pr[(x,y) \in \bigcup_{i \in [t]} (S_i \times B_i)] \enspace,  %
    \end{equation*}
    which is at least $\prox - \eps$, as required.
\end{proof}

We devote the rest of this section to proving \expref{Lemma}{lemma:key}.

\subsection{Proof of \expref{Lemma}{lemma:key}}

    Let $\blocklength = \abs{\Field}$, and define the binary matrix $A \in \{0,1\}^{\blocklength \times \blocklength}$
    where $A(x,y) = 1$ if $\colpoly(x,y) = \rowpoly(x,y)$ and $A(x,y) = 0$ otherwise.
    Note that by the assumption of \expref{Theorem}{thm:bivar-low-agreement}, we have
    \[
      \frac{\sum_{x,y \in [\blocklength]} A(x,y)}{\blocklength^2} = \prox\enspace,
      \]
      \ie, the matrix $A$ is $\prox$-dense.
\subsubsection{Step 1}


    In the first step we apply \expref{Theorem}{thm:KST} iteratively to show that there exists a collection of disjoint sets $S_1,\dots,S_{\numrect} \subseteq [\blocklength]$
    with $\abs{S_i} \geq {\degree}/{\eps}$ such that for most points $(x,y)$ it holds that if $A(x,y) = 1$,
%
%
%
    then $x \in \bigcup_i S_i$,             %
    and for each $i \in [\numrect]$ there exists $T_i \subseteq [\blocklength]$ of size $\abs{T_i} \geq {\degree}/{\eps}$ such that $A_{S_i \times T_i} \equiv 1$.

    For the rest of the proof we let $k = \ceil{\degree/\eps}$.

    \begin{claim}\label{claim:find all ones squares}
    Let $\blocklength,\degree \in \N$, $\prox > \eps >0$, and let $k \in \N$.
        Let $A \in \{0,1\}^{\blocklength \times \blocklength}$ be a $\prox$-dense matrix as above,
        and suppose that $\blocklength > 2k^2\left({1}/{\eps}\right)^{k+1}$.
        Then, there exist $\numrect \in \Naturals$ and two sequences $S_i \subseteq [\blocklength] ,T_i \subseteq [\blocklength]$ with $i=1, \dots, \numrect$
        satisfying the following conditions.
        \begin{enumerate}
%
%
            \item The sets $S_i$ are pairwise disjoint.                 %
            \item $\abs{S_i}=\abs{T_i} = k$.
            \item $A(x,y)=1$ for every $(x,y) \in S_i \times T_i$ and $i \in [\numrect]$.
%
            \item $\sum_{(x,y) \in  ([\blocklength] \setminus (\bigcup_i S_i)) \times [\blocklength]} A(x,y) < \eps \blocklength^2$. \label{item:leftover eps}      %
%
%
%
%
        \end{enumerate}
    \end{claim}

   \begin{proof}
        We will use \expref{Theorem}{thm:KST} to find a submatrix
        of $A$ of size $k \times k$ whose entries are all $1$s.
        By the choice of $k$ and the assumption that $\blocklength$ is sufficiently large
        we have that
        \[
          \left(\eps - \frac{k}{n}\right)^k = \eps^k\left(1 - \frac{k}{\eps n}\right)^k > \eps^k\left(1 - \frac{k^2}{\eps n}\right) >\eps^k/2 > \frac{k-1}{\eps n}\enspace,
        \]
        and hence
        \[
          \eps > \sqrt[k]{\frac{k-1}{\eps n}} + \frac{k}{n}\enspace.
        \]
          Hence, since $A$ is $\prox$-dense, we have
          \[
            \delta(A) \geq \prox \geq \eps >
          \sqrt[k]{\frac{k-1}{\blocklength}} +
          \frac{k}{\blocklength}\enspace.
        \]
        Therefore, by \expref{Theorem}{thm:KST} there exist
        $S_1 \subseteq [\blocklength], T_1 \subseteq [\blocklength]$
        each of size $\abs{S_1} = \abs{T_1} = k$ such that
        $A \vert_{S_1 \times T_1} \equiv 1$.

        Next, we remove the rows contained in $S_1$ from $A$,
        and apply the same argument again.
        Let $M_1 = [\blocklength] \setminus S_1$ and
        define $A_1$ to be the $(\blocklength-k) \times \blocklength$ submatrix of $A$
        whose rows are indexed by $M_1$.
        Note that if
        \[
          \sum_{x \in M_1, y \in [\blocklength]}A_1(x,y) > \eps \blocklength^2
        \]
        then $\delta(A_1) \geq {\eps \blocklength}/{\abs{M_1}}$,
        and thus we have
        \[
          \delta(A_1) \geq \frac{\eps \blocklength}{\blocklength-k} >
        \eps > \sqrt[k]{\frac{k-1}{\abs{M_1}}} +
        \frac{k}{\blocklength}\enspace.
      \]
      Therefore, we can apply \expref{Theorem}{thm:KST} again, and
      find $S_2 \subseteq M_1$ and $T_2 \subseteq [\blocklength]$ of
      size $\abs{S_2} = \abs{T_2} = k$ such that
      $A \vert_{S_2 \times T_2} \equiv 1$.

        We repeat the same argument again,
        for each $i \geq 2$ defining the
        the subset $M_i = M_{i-1} \setminus S_{i-1}$,
        and letting $A_i = A_{M_i \times [\blocklength]}$.
        Note that if
        \[
          \sum_{x \in M_i, y \in [\blocklength]} A(x,y) \geq \eps \blocklength^2
        \]
        then $\abs{M_i} \geq \eps \blocklength$,
        and
        \[
          \delta(A_i) \geq \frac{\eps \blocklength}{\abs{M_i}} \geq \eps > \sqrt[k]{\frac{k-1}{\abs{M_i}}} + \frac{k}{\blocklength}\enspace.
        \]
        Therefore, by \expref{Theorem}{thm:KST} there exist $S_i \subseteq M_i$ and $T_i \subseteq [\blocklength]$
        of size $\abs{S_i} = \abs{T_i} = k$ such that  $A \vert_{S_i \times T_i} \equiv 1$.

        We stop the process after $\numrect$ iterations when
        \[
          \sum_{x \in M_{\numrect}, y \in [\blocklength]} A(x,y) < \eps \blocklength^2\enspace.
          \]
          By the %
          definition
%
%
of the sets $S_i$ and $T_i$,                    %
this gives us the subsets with the desired properties.
    \end{proof}

    By the assumption $\abs{\Field} = n > \exp(\Omega(({\degree}/{\eps}) \log({1}/{\eps})))$
    in \expref{Theorem}{thm:bivar-low-agreement} and the choice of $k = \ceil{\degree/\eps}$, we have $\blocklength > 2k^2 \left({1}/{\eps}\right)^{k+1}$.
    Therefore, we can apply \expref{Claim}{claim:find all ones squares} on
 $A$ to get
%
the sets $S_i$ and $T_i$  %
    as in the claim.

 \subsubsection{Step 2}

    Next, we find subsets $B_i$ such that  $T_i \subseteq B_i \subseteq [\blocklength]$ for %
%
the sets $T_i$                  %
from the previous step, so that most points $(x,y)$ satisfying $A(x,y) = 1$ are contained in $\cup_{i=1}^{\numrect} (S_i \times B_i)$.
    \begin{claim}\label{claim:Si-Ti box}
        Let $\{S_i\}_{i=1}^{\numrect}$ be the sets from \expref{Claim}{claim:find all ones squares}.
        For each $i \in [\numrect]$ define
        \[
          B_i = \{ y_0 \in [\blocklength] : \sum_{x \in S_i} A(x,y_0) \geq \degree+1\}\enspace.
          \]
        Then
        \begin{enumerate}
        \item
            $\sum_{i \in [\numrect]} \sum_{\substack{x \in S_i \\ y \in [\blocklength] \setminus B_i}} A(x,y) \leq \eps \blocklength^2$.
        \item
            $A(x,y)=1$ for every $(x,y) \in S_i \times B_i$ and $i \in [\numrect]$.
    \end{enumerate}
    \end{claim}

    \begin{proof}
        The first item is by the choice of $k \geq \degree/\eps$.
        In each $i \in [\numrect]$ and $y \in [\blocklength] \setminus B_i$ it holds that
        less than $\eps$ fraction of the entries are ones,
        and hence the total number of ones in all $i \in [\numrect]$ and $y \in [\blocklength] \setminus B_i$
        is less that $\eps \blocklength^2$.
        Formally, we have
        \begin{equation*}
        \sum_{i \in [\numrect]} \sum_{\substack{x \in S_i \\ y \in [\blocklength] \setminus B_i}} A(x,y) \leq \sum_{i \in [\numrect]} \sum_{y \in [\blocklength] \setminus B_i} r  \leq u \cdot \blocklength \cdot \degree \leq \eps \blocklength^2 \enspace,
        \end{equation*}
        where the last inequality uses the fact that $u \leq \blocklength/k$, and $k \geq \degree/\eps$.

        To prove the second item, we use \expref{Corollary}{cor:completing boxes}.
        Suppose that $A(x_0,y_0)=0$ for some $x_0 \in S_i$ and $y_0 \in B_i$.
        By the assumption on $B_i$, it holds that $\abs{\{x \in S_i : A(x,y_0) = 1\}} \geq \degree+1$.
        Let $S = \{x_0\} \cup \{x \in S_i : A(x,y_0) = 1\}$,
        and let $T = \{y_0\} \cup T_i$, so that $A(x,y) = 1$ for all
        $(x,y) \in S \times T \setminus \{(x_0,y_0)\}$.
        Recall that, by definition of $A$,
        $\rowpoly(x,y) = \colpoly(x,y)$ for all such $(x,y)$,
        and hence, by \expref{Corollary}{cor:completing boxes}
        we also have $\rowpoly(x_0,y_0) = \colpoly(x_0,y_0)$, and thus $A(x_0,y_0) = 1$.
    \end{proof}

    \medskip\noindent
    Note that the ones not covered by $\bigcup_{i} (S_i \times B_i)$ are the
    $\leq \eps \blocklength^2$ ones omitted in \expref{Claim}{claim:find all ones squares}
    and the $\leq \eps \blocklength^2$ ones disregarded in the proof of \expref{Claim}{claim:Si-Ti box} above.
    Let us also disregard all
%
the $S_i$ and the $B_i$     %
such that $\abs{B_i} \leq \eps \blocklength$, and consider only the remaining subsets.
%
%
%
%
Note that the union of all
the   %
$S_i \times B_i$             %
with $\abs{B_i} \leq \eps \blocklength$ can %
contain     %
at most $\eps \blocklength^2$ ones.
    Redefining $\numrect$ to be the number of remaining sets, we get two collections of subsets $\{S_i \subseteq [\blocklength], B_i  \subseteq [\blocklength]\}_{i=1}^{\numrect}$ such that:
    \begin{enumerate}
%
%
        \item The sets $S_i$ are pairwise disjoint.                         %
        \item $\abs{B_i} > \eps n$ for all $i \in [\numrect]$.
        \item $\abs{ \bigcup_{i=1}^{\numrect} (S_i \times B_i)} \geq (\prox - 3\eps)\blocklength^2$.
        \item $A(x,y) = 1$ for every $(x,y) \in {S_i \times B_i}$ and $i \in [\numrect]$.
    \end{enumerate}
    In particular, by \expref{Lemma}{lemma:interpolation} for each $i=1, \dots, \numrect$
    there is a polynomial $P_i$ of degree $(\degree,\degree)$ such that
    $\rowpoly(x,y) = \colpoly(x,y) = P_i(x,y)$ for all $(x,y) \in S_i \times B_i$.

   \subsubsection{Step 3}
    Next, we observe that if two sets $B_i,B_j$ from the previous step have large intersection,
    then the corresponding polynomials $P_i$ and $P_j$ are equal.

    \begin{claim}\label{claim:Pi=Pj}
        Suppose that $\abs{B_i \cap B_j} \geq \degree+1$ for some $i \neq j \in [\numrect]$.
        Then $P_i = P_j$ and $B_i = B_j$.
    \end{claim}
    \begin{proof}
        Denote $B = B_i \cap B_j$.
        Note that, for each $y \in B$,
        $P_i(x,y) = \colpoly(x,y)$ for all $|S_i| = k > \degree+1$ values of $x \in S_i$,
        and hence $P_i(x,y) = \colpoly(x,y)$ for all $x \in [\blocklength]$.
        In particular, $P_i(x,y) = \colpoly(x,y)$ for all $(x,y) \in S_j \times B$.
        Therefore, ${P_i} \vert_{S_j \times B}  \equiv {P_j} \vert_{S_j \times B}$,
        and thus $P_i \equiv P_j$ by \expref{Lemma}{lemma:zero poly}.
        Applying \expref{Corollary}{cor:completing boxes}, we conclude that $P_i(x,y) = P_j(x,y) = \colpoly(x,y) = \rowpoly(x,y)$
        for all $(x,y) \in (S_i \cup S_j) \times (B_i \cup B_j)$.
        This implies that $B_i = B_j$, as required.
    \end{proof}

    \subsubsection{Completing the proof}

    In the last step we will show that there is a short list of $t \leq {2}/{\eps}$
    polynomials $Q_1,\dots, Q_t$ such that each of the
%
$P_i$ is in fact equal to one of the $Q_j$.   %
%
 Indeed, denote the number of distinct sets $B_i$ by $t$. %
    By \expref{Claim}{claim:Pi=Pj}, if $B_i \neq B_j$ then $|B_i \cap B_j| \leq \degree$,
    and thus by the inclusion-exclusion principle we have
    \begin{equation*}
      \blocklength \geq \abs{ \bigcup_{i=1}^t B_i } \geq \sum_{i=1}^t \abs{B_i} - \sum_{i \neq j} \abs{B_i \cap B_j}
      \geq t \cdot \eps \blocklength - \binom{t}{2} \degree \enspace,
    \end{equation*}
    where in the last inequality we used the bound $\abs{B_i} > \eps \blocklength$ for all $i$. The original $\numrect$ rectangles are disjoint $k \times k$ squares, so it must hold that $\numrect k \leq n$. Since $k \geq {\degree}/{\eps}$, this implies that $\numrect \degree \leq \eps n$. However, it also must hold that $t \leq \numrect$, since merging rectangles can only decrease the total number of rectangles. Therefore, $t \degree \leq \eps \blocklength$. Combining this with the first inequality, we get that
    \begin{equation*}
    \blocklength \geq t \cdot \eps \blocklength - {t \choose 2} \degree \geq t \left(\eps \blocklength - \frac{t \degree}{2} \right) \geq \left(\eps \blocklength - \frac{\eps \blocklength}{2}\right) \implies \blocklength \geq t \cdot \frac{\eps \blocklength}{2} \implies \frac{2}{\eps} \geq t \enspace.
    \end{equation*}
    Therefore $t \leq {2}/{\eps}$, as required.


%
%
\subsection{Proof of \expref{Proposition}{prop:list size lb}}
\label{sec:prop-proof}

We prove that in \expref{Theorem}{thm:bivar-low-agreement} it is \emph{necessary} to have a list of ${(\prox/\eps)}/{\polylog(\prox/\eps)}$ polynomials in order to explain all but $\eps$ of the agreements of $\rowpoly$ and $\colpoly$.

\medskip

The idea of the proof is the following.
Fix $N$ elements of $\Field$ arbitrarily, and (for simplicity of notation) identify these elements with $1,\dots,N$.
We define a infinite sequence of positive numbers $(a_i)_{i \geq 1}$ such that the series $\sum_{i} a_i$ converges,
and we choose a normalization constant $C$ so that
\begin{inparaenum}[(i)]
\item $\sum_{i=1}^N (C \cdot a_i)^2 = \delta$, and
\item for all $\eps > 0$ and $t = {(\prox/\eps)}/{\polylog(\prox/\eps)}$ it holds that $\sum_{i=1}^t (C \cdot a_i)^2 < \delta-\eps$.
\end{inparaenum}
Choose a collection of $N$ disjoint sets $A_i \subseteq \Field$ of size $\abs{A_i} = C \cdot a_i \cdot \abs{\Field}$.\footnote{We assume for simplicity that $C \cdot a_i \cdot \abs{\Field}$ are all integers.}
Then, we define row polynomials $\rowpoly(x,y)$ that are constant on every row $y$, and column polynomials $\colpoly(x,y)$ that are constant on every column $x$,
so that $\rowpoly(x,y) = \colpoly(x,y) = i$ if and only if $(x,y) \in A_i \times A_i$.
That is $\rowpoly(x,y) \equiv i$ for all $y \in A_i$, and $\colpoly(x,y) = i$ for all $x \in A_i$.
In particular, we have $\rowpoly(x,y) = \colpoly(x,y)$ if and only if $(x,y) \in \bigcup_{i} (A_i \times A_i)$,
and hence
\[
  \Pr_{(x,y) \in \Field^2}[\rowpoly(x,y) = \colpoly(x,y)] = \sum_{i=1}^N (C \cdot a_i)^2\enspace.
  \]
  \expref{Proposition}{prop:list size lb} follows rather immediately
  by the convergence properties of the partial sums
  $\sum_{i=1}^N a_i^2$ and the definition of $\rowpoly$ and
  $\colpoly$.

%
\medskip

We now proceed to the formal proof. We start by examining the two sequences $\{c_1(N)\}_{N \in \Naturals}$ and $\{c_2(N)\}_{N \in \Naturals}$ defined as follows. For all $i \geq 1$ let
\[
  a_i = \frac{1}{i (\log_2(i+1))^2}\enspace,
\]
and let $c_1(N) = \sum_{i = 1}^N a_i$ and
$c_2(N) = \sum_{i = 1}^N a_i^2$. Note that both $c_1(N)$ and $c_2(N)$
are non-decreasing and uniformly bounded%
\footnote{The partial sum $\sum_{i=1}^{N} a_i$ can be upper bounded by
  \[
    a_1 +a_2 + \sum_{i=3}^{N} \frac{1}{i (\log_2(i))^2} \leq a_1 + a_2
  + \int_{2}^{N} \frac{1}{x \log_2^2(x)} dx = a_1 + a_2 + -
  \frac{\ln^2(2)}{\ln(x)} \Big\vert_2^N = a_1 + a_2 + -
  \frac{\ln^2(2)}{\ln(N)} + \ln(2) \leq a_1 + a_2 + \ln(2),
\]
and the
  partial sum $\sum_{i=1}^{N} a_i^2$ is upper bounded by
  $\sum_{i=1}^{N} a_i$.}  sequences, and hence both converge as
$N \to \infty$.  We first show that
\[
  1 \geq \frac{c_2(N)}{c_1(N)^2} > \frac{1}{3}
\]
for every $N$. The ratio is clearly at most $1$, as
\[
  c_1(N)^2 = \left(\sum_{i = 1}^N a_i\right)^2 \geq \sum_{i = 1}^N a_i^2 = c_2(N)
\]
for every $N$. To show the other direction, we can compute that
$c_2(N) \geq 1$ and $c_1(N) \leq 1.63$ for all $N \geq 1$. It follows
that
\[
  \frac{c_2(N)}{c_1(N)^2} > \frac{1}{1.63^2} > \frac{1}{3}
\]
for every $N$.

Let $\delta, \eps$ be such that ${1}/{12} > \delta \geq \eps > 0$, and set $N = \ceil{{\log(1/\eps)}/{\eps}}$. Redefine $c_1 \DefineEqual c_1(N)$ and $c_2  \DefineEqual c_2(N)$. {}From the previous calculations, we have that
\[
  \prox < \frac{1}{12} < \frac{c_2}{4 c_1^2} < 1\enspace.
\]
Let $\Field$ be a field of
%
order  %
\[
  \abs{\Field} > \frac{3 c_2 \degree}{a_N^2\prox} = \frac{3 c_2 \degree N^2 \log^4(N+1)}{\prox} \geq  \frac{3 c_2 \log^6(1/\eps) \cdot \degree}{\prox \cdot \eps^2}\enspace.
\]
In particular, we have
\[
  \abs{\Field} > \frac{3 c_2 \degree}{a_N^2\prox} \geq \frac{3 c_2 \degree}{a_i^2\prox}
\]
for all $i \in [N]$.

We now show that there are disjoint subsets $A_1, \dots, A_N \subseteq \Field$ of size $\abs{A_i} = \ceil{\sqrt{\prox/c_2} \cdot a_i \abs{\Field}}$. We have that
\begin{equation*}
        \sum_{i=1}^N \abs{A_i}
        = \sum_{i=1}^N  \ceil{\sqrt{\prox/c_2} \cdot a_i \abs{\Field}}
        \leq \frac{4}{3} \sqrt{\prox/c_2} \sum_{i=1}^N  a_i \abs{\Field}
        \enspace,
    \end{equation*}
    where the last inequality holds since
    \[
      \sqrt{\prox/c_2} \cdot a_i \abs{\Field} \geq \frac{3}{a_i} \cdot \sqrt{\frac{c_2}{\prox}} \geq 3
    \]
    for every $i \in [N]$, by our assumption on $\abs{\Field}$. Since $\prox < c_2 / 4c_1^2$, it follows that $\sum_{i=1}^N \abs{A_i} < ({2}/{3})\abs{\Field}$, and hence
%
%
%
such sets $A_i$ exist.                  %

We now construct $\rowpoly(x,y)$ and $\colpoly(x,y)$ that satisfy the conditions in the proposition.
Fix $N$ elements of $\Field$ arbitrarily, and (for simplicity of notation) identify these elements with $1,\dots,N$.
For every $y \in A_i$, set the row polynomial $\rowpoly(x,y)$ equal to $i$, and for every $x \in A_i$,
    set the column polynomial $\colpoly(x,y)$ equal to $i$.
    For all $y \notin \bigcup A_i$ let $\rowpoly(x,y) = N+1$,
    and for all $x \notin \bigcup A_i$ let $\colpoly(x,y) = N+2$.
    Note that $\rowpoly(x,y) = \colpoly(x,y)$ if and only if $x,y \in A_i$ for some $i$.
    In particular,
    \begin{equation*}
        \Pr_{x,y}[\rowpoly(x,y) = \colpoly(x,y)] = \sum_{i=1}^N \left( \frac{\abs{A_i}}{\abs{\Field}} \right)^2
        = \sum_{i=1}^N \left(\frac{\ceil{\sqrt{\prox/c_2} \cdot a_i \abs{\Field}}}{\abs{\Field}} \right)^2
        \enspace.
    \end{equation*}
    The later expression is at least
    \[
      \sum_{i} \left(\sqrt{\frac{\prox}{c_2}} \cdot a_i\right)^2 = \sum_{i} \left(\frac{\prox}{c_2}\right) a_i^2 = \prox
    \]
    and at most
    \[
      \sum_{i} \left(\frac{4}{3} \sqrt{\frac{\prox}{c_2}} \cdot a_i\right)^2  < 2\sum_{i} \left(\frac{\prox}{c_2}\right) a_i^2  = 2\prox\enspace,
      \]
    and hence $\rowpoly$ and $\colpoly$ satisfy the requirement that
    \begin{equation*}
        \Pr_{x,y}[\rowpoly(x,y) = \colpoly(x,y)] = \sum_{i=1}^N \left( \frac{\abs{A_i}}{\abs{\Field}} \right)^2
        \in [\prox, 2\prox]
        \enspace.
    \end{equation*}
    By the %
%
%
%
%
%
%
%
%
    Polynomial Identity Lemma,\footnote{This is also known as Schwartz-Zippel Lemma,                %
      We refer to \cite{arvind} Section 3.1 for discussion about the history of the lemma.}           %
    any non-constant polynomial of individual degree $\degree$ (and therefore total degree $2 \degree$) can be equal to $i$
    on at most $2\degree \cdot \abs{A_i}$ points in each $A_i \times A_i$ box, and hence
    agrees with both $\rowpoly$ and $\colpoly$ on at most $2\degree \sum \abs{A_i} \leq 2\degree \abs{\Field}$
    points in total. Since
    \[
      \abs{\Field} > \frac{3 c_2 \degree}{a_i^2\prox}
    \]
    for every $i$, it follows that $\abs{A_i}^2 > 2\degree \abs{\Field}$ for every $i \in [N]$,
    and therefore the $t$ polynomials $Q_i$ that maximize the probability
    $\Pr_{x,y}[\exists i \in [t] \st \rowpoly(x,y) = \colpoly(x,y) = Q_i(x,y)]$
   must be the constant polynomials $Q_1 \equiv 1,\dots,Q_t \equiv t$.

   Since
   \[
     \Pr_{x,y}[\rowpoly(x,y) = \colpoly(x,y) = Q_i(x,y)] = \left(\frac{\abs{A_i}}{\abs{\Field}}\right)^2 \geq \frac{\prox}{c_2} a_i^2
     \]
    for each $i \in [N]$, we have
    \begin{eqnarray*}
        \Pr_{x,y}[\exists i \in [t] \st \rowpoly(x,y) = \colpoly(x,y) = Q_i(x,y)]
        &=& \Pr_{x,y}[\rowpoly(x,y) = \colpoly(x,y)] - \sum_{i=t+1}^N \left(\frac{\abs{A_i}}{\abs{\Field}}\right)^2 \\
        &\leq& \Pr_{x,y}[\rowpoly(x,y) = \colpoly(x,y)] - \sum_{i=t+1}^N \frac{\prox}{c_2} a_i^2
        \enspace.
    \end{eqnarray*}
    Note that
    \[
      \sum_{i=t+1}^\infty a_i^2 = \Theta \left( \frac{1}{t \cdot \polylog(t)} \right)\enspace,
    \]
    and since $N > {\log(1/\eps)}/{\eps}$
    it holds that
    \[
      \sum_{i=N}^{\infty} \frac{\prox}{c_2} a_i^2 < \frac{\prox}{c_2}\eps < \eps\enspace.
    \]
    Thus,
    \[
      \sum_{i=t+1}^N \frac{\prox}{c_2} a_i^2 \geq \Omega \left( \frac{\prox}{t \cdot \polylog(t)} \right) - \eps\enspace.
    \]
    This implies that
    \begin{equation*}
        \Pr_{x,y}[\exists i \in [t] \st \rowpoly(x,y) = \colpoly(x,y) = Q_i(x,y)] < \Pr_{x,y}[\rowpoly(x,y) = \colpoly(x,y)] - O\left(\frac{\prox}{t \cdot \polylog(t)}\right)+ \eps
        \enspace.
    \end{equation*}
    Therefore, if the left hand side of the above is at least $\Pr_{x,y}[\rowpoly(x,y) = \colpoly(x,y)] - \eps $, then we have that $t  > {(\prox/\eps)}/{\polylog(\prox/\eps)}$,
    as required.

%
%
%
\section{Preliminaries for \expref{Theorem}{thm:pat-robustness}}
\label{sec:prelim-pat}

%
%
\subsection{Linear codes}
\label{sec:linear-codes}

A linear code $C \subseteq \Field^{\blocklength}$ is a linear subspace of the vector space $\Field^{\blocklength}$. Each codeword $\word$ in $C$ is a string of length $\blocklength$, which is the block length of the code. The dimension of the code $\dim(C)$ is the dimension of $C$ as a vector space in $\Field^n$. For any two words $\word$ and $\otherword$ in $\Field^{\blocklength}$, the Hamming distance between $\word$ and $\otherword$, denoted by $\hamdist(\word,\otherword)$, is the number of indices where $i$ where $\word_i \ne \otherword_i$. Formally, $\hamdist(\word,\otherword) = \abs{\{i \in [\blocklength] : \word_i \ne \otherword_i\}}$. The relative distance between $\word$ and $\otherword$ is $\reldist(\word,\otherword) = \hamdist(\word,\otherword)/\blocklength$, which is the fraction of points where $\word$ and $\otherword$ disagree. For any subset $S$ of $[\blocklength]$, we will define $\hamdist \vert_S(\word,\otherword)$ to be $\abs{\{i \in S : \word_i \ne \otherword_i\}}$, which is the Hamming distance between $\word$ and $\otherword$ on the subset $S$. Similarly, $\reldist \vert_S(\word,\otherword) = \hamdist \vert_S (\word,\otherword) / \abs{S}$. The distance $\codedist$ of a code $C$ is the minimum Hamming distance between any two distinct codewords of $C$, \ie, $d = d(C) = \min_{\word \neq \otherword \in C} \hamdist(\word,\otherword)$. For any $\word$ in $\Field^{\blocklength}$, the distance from $\word$ to $C$ is defined as $\hamdist(\word,C) =\min_{\otherword \in C} \hamdist(\word,\otherword)$, and the relative distance is defined similarly. For any subset $S \subseteq [\blocklength]$, the distance from $\word$ to $C$ on $S$ is $\hamdist \vert_S (\word, C) = \min_{\otherword \in C} \hamdist \vert_S(\word,\otherword)$. We will write $\reldist(\word)$ instead of $\reldist(\word,C)$ when the code is clear from the context.

Linear codes have a unique extension property.
\begin{claim}[Unique Extension]\label{claim:ext}
Let $C \subseteq \Field^{\blocklength}$ be a code with blocklength $\blocklength$ and distance $\codedist$.
Let $I$ be a subset of $[\blocklength]$ of size at least $\blocklength - \codedist + 1$. Let $C'$ be the restriction of the code $C$ to the subset $I$.
Then, for every codeword $\word \in C'$ there exists a unique $\otherword \in C$ such that $\otherword \vert_I = \word$.
\end{claim}
\begin{proof}
By definition, for every $\word$ in $C'$ there must exist at least one $\otherword$ in $C$ such that $\otherword \vert_I = \word$. Suppose there exists $\otherword_1$ and $\otherword_2$ such that $\otherword_1 \vert_I = \otherword_2 \vert_I = \word$. This implies that $\hamdist(\otherword_1, \otherword_2) \leq \blocklength - \abs{I} \leq \codedist - 1$, as $\otherword_1$ and $\otherword_2$ agree on $I$. Since $\otherword_1$ and $\otherword_2$ are codewords, $\hamdist(\otherword_1,\otherword_2) < \codedist$ if and only if $\otherword_1 = \otherword_2$. Therefore, the codeword $\word$ has a unique extension to $C$.
\end{proof}
%
%
\subsection{Tensor product codes}
\label{sec:tensor-product-codes}

%
%
%

For two linear codes
$C_1 \subseteq \Field^{\blocklength_1}$ and
$C_2 \subseteq \Field^{\blocklength_2}$,
the tensor product of $C_1$ and $C_2$, denoted by $C_1 \otimes C_2$,
is the linear code in
$\Field^{\blocklength_1\cdot\blocklength_2}$,
where every codeword
$\testerinput \in \Field^{\blocklength_1 \times \blocklength_2}$
is an $\blocklength_1 \times \blocklength_2$ matrix, whose each column
is a codeword of $C_1$ and each row is a codeword of $C_2$.
In particular, we have $\dim(C_1\otimes C_2)=\dim(C_1)\cdot\dim(C_2)$
and $\codedist(C_1\otimes C_2) = \codedist(C_1) \cdot \codedist(C_2)$.
%
%
%
%
We then extend this definition inductively to $m\ge 3$ codes
$C_i \subseteq \Field^{\blocklength_i}$, setting
$C_1\otimes\dots\otimes C_m:=
(C_1\otimes\dots\otimes C_{m-1})\otimes C_m$.  
So $C_1\otimes\dots\otimes C_m$ has
block length $\blocklength_1\cdots \blocklength_m$,
distance $\codedist(C_1)\cdots \codedist(C_m)$, and
dimension $\dim(C_1)\cdots\dim(C_m)$.

Furthermore, $\closestglobalword \in C_1\otimes\dots\otimes C_m$
if and only if it can be written as an $m$-dimensional 
$\blocklength_1 \times \blocklength_2 \times \dots \times \blocklength_m$
tensor, where the entries are values in $\Field$, and each line parallel
to the $i$-th axis is in $C_i$. It is easy to see that
$\closestglobalword$ is in $C_1\otimes\dots\otimes C_m$
if and only if for each $i$, the restriction of $\closestglobalword$ to
any $(m-1)$-dimensional axis-parallel hyperplane perpendicular to
the $i$-th axis is in
$C_1\otimes\cdots\otimes C_{i-1}\otimes C_{i+1}\otimes\cdots\otimes C_m$.

If $C_1=\dots=C_m=C$ then we call the
product $C_1\otimes\dots\otimes C_m$ the
\emph{$m$-th tensor power} of $C$ and denote it by $C^m$.
It is worth noting that the fractional distance of the code
$C^m$ is $(\codedist/\blocklength)^m$, so the fractional distance
of the code decays exponentially in $m$.
%

%
%
%
%
%
%
%
%
%
%
%
%
%
%
%
%
%
%
%
%
%
%
%
%
%
%
%
%
%
%

Tensor product codes have a unique extension property that will be used many times in the proof of \expref{Theorem}{thm:pat-robustness}.

\begin{claim}[Unique Extension for Tensor Product Codes]
\label{claim:tensorext}
Let $\{C_b\}_{b = 1}^{m}$ be codes with blocklength $\blocklength_b$ and distance $\codedist_b$. Let $I_b \subseteq [\blocklength_b]$ be a set of size at least $\blocklength_b - \codedist_b + 1$, and let $C'_b$ be the projection of $C_b$ to $I_b$. Then for every $\word \in C' =  C'_1 \otimes \dots \otimes C'_m$, there exists a unique $\otherword$ in $C = C_1 \otimes \dots \otimes C_m$ such that $\otherword \vert_{I_1 \times \dots \times I_m} = \word$.
\end{claim}

\begin{proof}
By \expref{Claim}{claim:ext}, for all $b \in[m]$ the projection map $\pi_b \colon C_b \to C'_b$ is bijective. We can extend $\pi_b$ to be a bijective map from the hybrid code $C'_1 \otimes \dots \otimes C'_{b-1} \otimes C_b \otimes \dots \otimes C_m$ to $C'_1 \otimes \dots \otimes C'_{b} \otimes C_{b+1} \otimes \dots \otimes C_m$. For any $\otherword$ in the first hybrid code, define $\pi_b(\otherword) = \otherword \vert_{I_1 \times \dots \times I_b \times \blocklength_{b+1} \times \dots \times \blocklength_m}$, which is the projection of $\otherword$ to $I_b$ along the $b$th axis, and the identity map everywhere else. Clearly, $\pi_b$ is still a bijection, and so the composition of maps $\pi = \pi_m \circ \pi_{m-1} \circ \dots \circ \pi_1$ is therefore a bijection from $C$ to $C'$, which proves the claim.
\end{proof}



%
%
\subsection{Locally testable codes and robust tests}
\label{sec:ltc}

A $q$-query test $\arbtester$ for a code $C \subseteq \Field^{\blocklength}$ is a probabilistic algorithm that, given oracle access to a word $\word \in \Field^{\blocklength}$, makes $q$ (non-adaptive) queries to $\word$ and then accepts or rejects. Informally, $C$ is locally testable if there is a test $\arbtester$ that accepts (with probability $1$) whenever $\word$ is in $C$, and rejects (say with probability at least $0.5$) when $\word$ is far from $C$.

The expected local view distance $\robust^{\arbtester}(\word)$ of $\arbtester$ on a word $\word$ is the average, over the local views of $\arbtester$, of the distance of $\word$ to an accepting view. Instead of analyzing the local testability of $C^m$, we will instead consider a stronger notion of local testability called robustness, that was introduced in \cite{BenSassonGHSV06} and adapted to locally  testable codes in \cite{Ben-SassonS06}. The test $\arbtester$ is \emph{$\alpha$-robust} if $\robust^{\arbtester}(\word) \geq \alpha \cdot \reldist(\word, C)$ for every word $\word \in \Field^{\blocklength}$. The ``strength'' of the test increases with $\alpha$, so the goal is to establish the largest $\alpha$ for which this inequality holds.

For formal definitions and more details on the subject see, \eg, \cite{Goldreich2010}.
%
%
\subsection{The axis-parallel hyperplane test}

Given a linear code $C$, we define below the \emph{axis-parallel hyperplane test} for the tensor code $C^m$.
An $(m-1)$-dimensional axis-parallel hyperplane $\plane \subseteq [\blocklength]^m$ is a set of the form $\plane = \{p \in [\blocklength]^m : p_b = i\}$ for some $b \in [m]$ and $i \in [\blocklength]$. Therefore, there are $\blocklength m$ such $(m-1)$-dimensional axis-parallel hyperplanes in total, and each such hyperplane can be specified by the pair $(b,i)$. We will use the notation $(b,i)$ to refer to the hyperplane $H = \{p \in [\blocklength]^m : p_b = i\}$.


\begin{definition}\label{def:tester}
Let $C$ be a linear code, and let $C^m$ be the
%
$m$-th tensor power of $C$.                                                     %
The axis-parallel hyperplane test $\tester$ for $C^m$ is the test that
given a word $\testerinput \in \Field^{\blocklength^{m}}$
samples a random $(m-1)$-dimensional axis-parallel hyperplane $\plane$
and checks if $\testerinput \vert_{\plane} \in C^{m-1}$.
\end{definition}

For $\testerinput \in \Field^{\blocklength^{m}}$ and an axis-parallel hyperplane $\plane$ in $[\blocklength]^m$, we define $\closestlocalword{\plane}$ to be the closest $C^{m-1}$ codeword to $\testerinput \vert_{\plane}$.
If this codeword is not unique, then we break ties by picking an arbitrary closest codeword.
Using this notation, the expected local view distance $\robust(\testerinput)$ can be expressed as
\begin{equation*}
\robust(\testerinput) = {\mathbb E}_{\plane}[\reldist \vert_{\plane}(\testerinput, C^{m-1})]
 = {\mathbb E}_{\plane}[\reldist \vert_{\plane}(\testerinput, \closestlocalword{\plane})]
%
\enspace,
\end{equation*}
where the expectation is taken over all axis-parallel hyperplanes $\plane$.

\begin{definition}\label{def:robust tester}
The test $\tester$ is \emph{$\alpha$-robust} if $\robust(\testerinput) \geq \alpha \cdot \reldist(\testerinput, C^{m})$ for every word $\testerinput \in \Field^{\blocklength^{m}}$, where $\reldist(\testerinput, C^{m})$ is the relative distance of the word $\testerinput$ to the code $C^{m}$, and $\robust(\testerinput)$ is the expected local distance of $\testerinput$.
\end{definition}

Note that robustness $\alpha$ for the test $\tester$ is at most $1$.
\begin{lemma}
The robustness of the axis-parallel hyperplane test $\tester$ is $\alpha \leq 1$.
\end{lemma}
\begin{proof}
Let $\closestglobalword$ be any $C^m$ codeword such that $\reldist(\testerinput) = \reldist(\testerinput, \closestglobalword)$. Then,
\begin{flalign*}
\reldist(\testerinput) = \reldist(\testerinput, \closestglobalword) = \frac{1}{\blocklength m} \sum_{\plane} \reldist \vert_{\plane} (\testerinput, \closestglobalword) \geq  \frac{1}{ \blocklength m} \sum_{\plane} \reldist \vert_{\plane} (\testerinput, \closestlocalword{\plane}) = \robust(\testerinput)
\end{flalign*}
since $\closestlocalword{\plane}$ is closer to $\testerinput \vert_{\plane}$ than $\closestglobalword \vert_{\plane}$, as $\closestglobalword \vert_{\plane} \in C^{m-1}$. Thus $\alpha \leq \robust(\testerinput) / \reldist(\testerinput) \leq 1$.
\end{proof}

In \expref{Section}{sec:soundness} we discuss the difference between robustness and soundness for the test $\tester$.



%
%
%
\section{Proof of \expref{Theorem}{thm:pat-robustness}}
\label{sec:proof1}

We prove \expref{Theorem}{thm:pat-robustness}. We start by recalling some notation. Let $C$ be a linear code with distance $\codedist$ and block length $\blocklength$ over $\Field$, and let $C^m$ be the
%
$m$-th tensor power of $C$,                                                             %
for some $m \geq 3$. Let $\testerinput$ be the input to the test $\tester$, which is an evaluation table of a function from $\mathbb [\blocklength]^m \to \Field$. Define $\closestlocalword{\plane}$ to be the closest $C^{m-1}$ word to $\testerinput \vert_{\plane}$, where ties are broken by picking an arbitrary closest codeword. We will view $\testerinput$ as fixed throughout the analysis.

Our goal is to show that $\robust(\testerinput) \geq \alpha \cdot \reldist(\testerinput)$, for
\[
  \alpha = \frac{1}{12} \left(\frac{\codedist}{\blocklength}\right)^m\enspace.
\]
In order to do that we assume that $\robust(\testerinput)$ is small, and use this to deduce that $\reldist(\testerinput)$ is also small. The assumption that $\robust(\testerinput)$ is small means that
%
%
the restrictions $\testerinput \vert_{\plane}$              %
are close on average to $\closestlocalword{\plane}$, and the proof uses this assumption to ``stitch'' together the
%
$\closestlocalword{\plane}$    %
into a global codeword $\closestglobalword$ that is close to $\testerinput$.

We begin by defining the inconsistency graph $G$. The graph $G$ has each hyperplane as a vertex, and has an edge between two hyperplanes $\plane$ and $\otherplane$ if they have nonzero intersection and their respective local codewords $\closestlocalword{\plane}$ and $\closestlocalword{\otherplane}$ are inconsistent, \ie, they disagree on some point $p$ in their intersection $\plane \cap \otherplane$.

%
\begin{definition}[Inconsistency Graph]\label{def:graph}
The inconsistency graph $G$ of the test $\tester$ is a graph where $V$ is the set of hyperplanes, and $E = \{(\plane, \otherplane) : \exists p \in \plane \cap \otherplane \st \closestlocalword{\plane}(p) \ne \closestlocalword{\otherplane}(p)\}$.
\end{definition}
%

The proof will be divided into several steps. First, we will show that if $G$ contains a large independent set, namely a large set of planes which are all consistent with each other, then there is a global codeword $\closestglobalword$ that stitches together all of the local codewords $\closestlocalword{\plane}$ for every $\plane$ in the independent set. Then, we will show that every edge in $G$ is adjacent to a vertex of (somewhat) large degree. This will imply that the set of vertices that have large degree is a vertex cover, and its complement is an independent set.
We will then
%
give an upper bound on  %
the number of vertices with large degree in terms of $\robust(\testerinput)$, and therefore if $\robust(\testerinput)$ is small then the vertex cover is small and hence the independent set is large.
Using these components we will conclude the proof.

%
%
\subsection{Step 1: the case of a large independent set}

We will show that if $G$ has a large independent set $I$, then there is an $\closestglobalword$ in $C^m$ that agrees with $\closestlocalword{\plane}$ on $\plane$ for every $\plane$ in $I$. In other words, $\closestglobalword$ is the codeword of $C^m$ that stitches together all of the
%
$\closestlocalword{\plane}$      %
in the independent set. The proof relies on \expref{Claim}{claim:tensorext}.
\begin{lemma}[Interpolation]\label{lemma:interp}
If $G$ has an independent set $I$ of size $\abs{I} > (m-1)(\blocklength-\codedist) + \blocklength$, then there exists $\closestglobalword$ in $C^m$ such that $\closestglobalword \vert_{\plane} = \closestlocalword{\plane}$ for every $\plane \in I$.
\end{lemma}
\noindent
Our proof of this lemma is similar to the proof of a different lemma in \cite{Ben-SassonS06}.
\begin{proof}
Define $I_b$ to be the set of $i \in [\blocklength]$ such that the plane $(b,i)$ is in $I$. Since $\abs{I} > (m-1)(\blocklength-\codedist) + \blocklength$, there must exist $b_1 \ne b_2$ such that $\abs{I_{b_1}}$ and $\abs{I_{b_2}}$ are at least $\blocklength - \codedist + 1$, as otherwise $\abs{I} = \sum_{b= 1}^m \abs{I_b} \leq (m-1)(\blocklength - \codedist) + \blocklength$. Without loss of generality assume $b_1 = 1$ and $b_2 = 2$. Let $S = I_1 \times I_2 \times [\blocklength]^{m-2}$ and let $g \colon S \to \Field$ be a matrix in $\Field^{S}$. Define $g(p) = \closestlocalword{\plane}(p)$ for every $p \in S$, where $\plane$ is some plane in $I_1 \cup I_2$ such that $p \in \plane$. Note that $g$ is well-defined since all the planes in $I$ are consistent with each other, as $I$ is an independent set.

We claim that $g \in C \vert_{I_1} \otimes C \vert_{I_2} \otimes C^{m-2}$. This is because for any $\plane \in I_1$ it holds that $g \vert_{\plane} \in C \vert_{I_2} \otimes C^{m-2}$, as $g \vert_{\plane} = \closestlocalword{\plane}$ except that the second axis is now restricted to $I_2$.
This means that for every axis $b \ne 1,2$ and for every line $\ell_b$ parallel to the $b$-th axis it holds that $g \vert_{\ell_b} \in C$.
Also, for every line $\ell_2$ parallel to the second axis we have that $g \vert_{\ell_2} \in C \vert_{I_2}$, because we took a $C^{m-1}$ codeword and restricted it to the subset $I_2$.
However, by symmetry we can repeat the same argument, swapping axis $1$ and axis $2$, and hence for every line $\ell_1$ parallel to the first axis it must hold that $g_{\ell_1} \in C \vert_{I_1}$.
Thus, $g \in C \vert_{I_1} \otimes C \vert_{I_2} \otimes C^{m-2}$. Since $\abs{I_1}$ and $\abs{I_2}$ are at least $\blocklength - \codedist + 1$, we can apply \expref{Claim}{claim:tensorext} to the code $C \vert_{I_1} \otimes C \vert_{I_2} \otimes C^{m-2}$ to extend $g$ to a unique codeword $\closestglobalword \in C^m$.


We still need to show that $\closestglobalword \vert_{\plane} = \closestlocalword{\plane}$ for every $\plane \in I$. By definition of $C^m$ we have $\closestglobalword \vert_{\plane} \in C^{m-1}$. There are three cases. If $\plane \in I_1$, then $\closestglobalword$ agrees with $\closestlocalword{\plane}$ on a subset of $\plane$ of size $I_2 \times [\blocklength]^{m-2}$, because $\closestlocalword{\plane} \vert_{I_2 \times [\blocklength]^{m-2}} = g \vert_{I_2 \times [\blocklength]^{m-2}} = \closestglobalword \vert_{I_2 \times [\blocklength]^{m-2}}$. Similarly, if $\plane \in I_2$, then $\closestglobalword$ agrees with $\closestlocalword{\plane}$ on a subset of size $I_1 \times [\blocklength]^{m-2}$, and if $\plane \in I \setminus (I_1 \cup I_2)$, then $\closestglobalword$ agrees with $\closestlocalword{\plane}$ on a subset of size $I_1 \times I_2 \times [\blocklength]^{m-3}$.
In all $3$ of the cases, since $\abs{I_1}$ and $\abs{I_2}$ are at least $\blocklength - \codedist + 1$, by
\expref{Claim}{claim:tensorext} there is a unique codeword $\word \in C^{m-1}$ that equals $\closestglobalword \vert_{\plane}$ (or $\closestlocalword{\plane}$) on that subset of $\plane$. But $\closestglobalword \vert_{\plane}$ is in $C^{m-1}$, so by the uniqueness of the extension it follows that $\closestglobalword \vert_{\plane} = \closestlocalword{\plane}$.
\end{proof}

%
%
\subsection{Step 2: the structure of \texorpdfstring{$G$}{G}}

We will now show that every edge $(\plane, \otherplane)$ in $G$ is adjacent to a vertex of large degree. The proof uses the structure of $C^m$ to show that if two planes disagree on a point, they must disagree on many points, and these points have a certain structure. Using the structure of these points, we find $(m-2)\codedist$ planes that intersect $\plane \cap \otherplane$ on at least one point
%
where   %
$\closestlocalword{\plane}$ and $\closestlocalword{\otherplane}$ disagree, and therefore each of these new planes must be adjacent to at least one of $\plane$ and $\otherplane$.
\begin{lemma}\label{lemma:vc}
If $(\plane, \otherplane) \in E$, then $\deg(\plane) + \deg(\otherplane) \geq (m-2) \codedist$.
\end{lemma}
\noindent A similar lemma appears in \cite{Ben-SassonS06}, but the graph they consider is different from ours.
\begin{proof}
Without loss of generality assume that $\plane = (1,i)$ and $\otherplane = (2,j)$. Fix $k \in \{3, \dots, m\}$. Let $I_k$ be the
%
%
set of indices $\ell \in [\blocklength]$ such that                          %
%
the plane $(k,\ell)$ is not adjacent to both $\plane$ and $\otherplane$. Suppose $\abs{I_k} \geq \blocklength - \codedist + 1$. Then $\closestlocalword{\plane} \vert_{I_k \times [\blocklength]^{m-3}} = \closestlocalword{\otherplane} \vert_{I_k \times [\blocklength]^{m-3}}$. Since $\abs{I_k} \geq \blocklength - \codedist + 1$, by \expref{Claim}{claim:tensorext} $\closestlocalword{\plane} \vert_{I_k \times [\blocklength]^{m-3}}$ can be extended to a unique $\word \in C^{m-2}$, and so $\word = \closestlocalword{\plane} \vert_{\plane \cap \otherplane}$. Similarly, $\closestlocalword{\otherplane} \vert_{I_k \times [\blocklength]^{m-3}}$ can be extended to a unique $\otherword \in C^{m-2}$, and so $\otherword = \closestlocalword{\otherplane} \vert_{\plane \cap \otherplane}$. However, since both $ \closestlocalword{\plane} \vert_{\plane \cap \otherplane}$ and $ \closestlocalword{\otherplane} \vert_{\plane \cap \otherplane}$ agree on $I_k \times [\blocklength]^{m-3}$, the uniqueness of the extension implies that they are equal, contradicting the fact that $(\plane, \otherplane)$ is an edge in the graph. Therefore, $\abs{I_k} \leq \blocklength - \codedist$ for every $k$. This means that for a fixed $k$, there are at least $\codedist$ planes $(k,\ell)$ such that $\closestlocalword{\plane}$ and $\closestlocalword{\otherplane}$ disagree on the intersection of all $3$ planes. Since $\closestlocalword{\plane}$ and $\closestlocalword{\otherplane}$ disagree, $\closestlocalword{(k,\ell)}$ can agree with at most one of them, so at least one of $(\plane, (k,\ell))$ and $(\otherplane, (k,\ell))$ is an edge. This holds for at least $\codedist$ planes for every $k$, which is a total of $(m-2)\codedist$ planes. Therefore, $\deg(\plane) + \deg(\otherplane) \geq (m-2)\codedist$.    %
\end{proof}
%
Thus, for every edge $(\plane, \otherplane)$ one of $\plane$ and $\otherplane$ has degree $\geq (m-2)\codedist/2$, so we deduce the following corollary.

\begin{corollary}[Vertex Cover]\label{cor:vc}
The set $L$ of vertices %
of   %
degree $\geq (m-2)\codedist/2$ is a vertex cover.
\end{corollary}

%
%
\subsection{Step 3: relating the expected local distance to the vertex cover}

We now relate the set of vertices of large degree to the expected local view distance of the test $\tester$. The main idea is to put the expression for $\robust(\testerinput)$ into a particular form, and then apply the triangle inequality to express $\robust(\testerinput)$ as a sum over edges in the graph. Using a simple relation between $\abs{L}$ and $\abs{E}$, the lemma follows.

\begin{lemma}\label{lemma:robust}
Let $L$ be the set of vertices %
of  %
degree $\geq (m-2)\codedist/2$ as in \expref{Corollary}{cor:vc}. Then
\[
  \robust(\testerinput) \geq \frac{m-2}{4(m-1)} \frac{\codedist^{m-1}}{\blocklength^{m-1}} \frac{\abs{L}}{\blocklength m}\enspace.
\]
\end{lemma}
\begin{proof}
  By definition,
  \[
    \robust(\testerinput) = \frac{1}{\blocklength^m m} \sum_{\plane} \hamdist(\testerinput \vert_{\plane}, \closestlocalword{\plane})\enspace.
  \]
  For any $\plane = (b,i)$,
\begin{equation*}
\hamdist(\testerinput \vert_{\plane}, \closestlocalword{\plane}) = \frac{1}{m-1} \sum_{c \in [m] \setminus \{b\}} \sum_{j \in [n]} \hamdist\vert_{\plane \cap (c,j)}(\testerinput, \closestlocalword{\plane}) = \frac{1}{m-1} \sum_{\otherplane : \plane \cap \otherplane \ne \emptyset} \hamdist\vert_{\plane \cap \otherplane}(\testerinput, \closestlocalword{\plane}) \enspace.
\end{equation*}
This is because for any point $p \in \plane$ and for any axis $c \ne b$, the point $p$ is in the intersection $\plane \cap (c,j)$ for exactly one $j$. Therefore,
\begin{equation*}
\robust(\testerinput) = \frac{1}{\blocklength^m m} \sum_\plane \hamdist(\testerinput \vert_\plane, \closestlocalword{\plane}) = \frac{1}{\blocklength^m m} \sum_\plane \frac{1}{m-1} \sum_{\otherplane : \plane \cap \otherplane \ne \emptyset} \hamdist \vert_{\plane \cap \otherplane} (\testerinput, \closestlocalword{\plane})
\end{equation*}
Every pair $(\plane, \otherplane)$ with $\plane \cap \otherplane \ne \emptyset$ appears exactly twice in the sum, contributing $\hamdist \vert_{\plane \cap \otherplane} (\testerinput, \closestlocalword{\plane})$ and $\hamdist \vert_{\plane \cap \otherplane} (\testerinput, \closestlocalword{\otherplane})$ to the sum. Therefore,
\begin{flalign*}
&\robust(\testerinput) = \frac{1}{\blocklength^m m(m-1)} \sum_{(\plane, \otherplane) : \plane\cap \otherplane \ne \emptyset} \hamdist \vert_{\plane \cap \otherplane} (\testerinput, \closestlocalword{\plane}) + \hamdist \vert_{\plane \cap \otherplane} (\testerinput, \closestlocalword{\otherplane})\\
& \geq  \frac{1}{\blocklength^m m(m-1)} \sum_{(\plane, \otherplane) : \plane\cap \otherplane \ne \emptyset} \hamdist \vert_{\plane \cap \otherplane} (\closestlocalword{\plane}, \closestlocalword{\otherplane}) = \frac{1}{\blocklength^m m(m-1)} \sum_{(\plane, \otherplane)\in E} \hamdist \vert_{\plane \cap \otherplane} (\closestlocalword{\plane}, \closestlocalword{\otherplane}) \enspace.
\end{flalign*}
as $(\plane, \otherplane) \notin E \implies \hamdist \vert_{\plane \cap \otherplane} (\closestlocalword{\plane}, \closestlocalword{\otherplane}) = 0$ by definition. Fix $(\plane, \otherplane) \in E$. The local codewords $\closestlocalword{\plane}$ and $\closestlocalword{\otherplane}$ are both in $C^{m-1}$, so $\closestlocalword{\plane} \vert_{\plane \cap \otherplane}$ and $\closestlocalword{\otherplane} \vert_{\plane \cap \otherplane}$ are both $C^{m-2}$ codewords. In particular, since $\hamdist \vert_{\plane \cap \otherplane} (\closestlocalword{\plane}, \closestlocalword{\otherplane}) > 0$, they are {\it distinct} codewords, and so $\hamdist \vert_{\plane \cap \otherplane} (\closestlocalword{\plane}, \closestlocalword{\otherplane}) \geq \codedist^{m-2}$. Therefore,
\begin{flalign*}
\robust(\testerinput) \geq \frac{1}{\blocklength^m m(m-1)} \sum_{(\plane, \otherplane)\in E} \hamdist \vert_{\plane \cap \otherplane} (\closestlocalword{\plane}, \closestlocalword{\otherplane}) \geq \frac{ \abs{E} \codedist^{m-2}}{\blocklength^m m(m-1)} \enspace.
\end{flalign*}
Since $L$ is the set of vertices of degree $\geq (m-2)\codedist/2$,
\begin{equation*}
2\abs{E} = \sum_{\plane} \deg(\plane) \geq \sum_{\plane\in L} \deg(\plane)\geq \abs{L} \frac{(m-2)\codedist}{2} \implies \abs{E} \geq \abs{L} \frac{(m-2)\codedist}{4} \enspace.
\end{equation*}

Thus,
\begin{equation*}
\robust(\testerinput) \geq \frac{ \abs{E} \codedist^{m-2}}{\blocklength^m m(m-1)} \geq  \frac{ (m-2)\abs{L} \codedist^{m-1}}{4\blocklength^m m(m-1)} = \frac{(m-2)}{4 (m-1)} \frac{\codedist^{m-1}}{\blocklength^{m-1}} \frac{\abs{L}}{\blocklength m} \enspace. \qedhere
\end{equation*}
\end{proof}

%
%
\subsection{Putting things together}

We are now ready to prove \expref{Theorem}{thm:pat-robustness}. The result follows from straightforward applications of the previous steps.
\begin{proof}[Proof of \expref{Theorem}{thm:pat-robustness}]
  Suppose first that
  \[
    \robust(\testerinput) \geq \frac{m-2}{4m} \frac{\codedist^m}{\blocklength^m}\enspace.
  \]
  Then we trivially have
  \[
    \robust(\testerinput) \geq \frac{m-2}{4m} \frac{\codedist^m}{\blocklength^m} \geq \reldist(\testerinput)\enspace,
  \]
  as $\reldist(\testerinput) \leq 1$.

  Now, suppose that
  \[
    \robust(\testerinput) < \frac{m-2}{4m} \frac{\codedist^m}{\blocklength^m}\enspace.
  \]
  Then by \expref{Lemma}{lemma:robust} we have
\begin{equation*}
\frac{m-2}{4m} \frac{\codedist^m}{\blocklength^m} > \robust(\testerinput) \geq  \frac{(m-2)}{4 (m-1)} \frac{\codedist^{m-1}}{\blocklength^{m-1}} \frac{\abs{L}}{\blocklength m} \enspace,
\end{equation*}
and so $\abs{L} < (m-1)\codedist$.
For every $\closestglobalword$ in $C^m$, using triangle inequality we have
\begin{equation*}
\reldist(\testerinput) \leq \reldist(\testerinput, \closestglobalword) = \frac{1}{\blocklength m} \sum_{\plane} \reldist \vert_{\plane} (\testerinput, \closestglobalword) \leq \frac{1}{\blocklength m} \sum_{\plane} \reldist \vert_{\plane}(\testerinput, \closestlocalword{\plane}) + \frac{1}{\blocklength m} \sum_{\plane} \reldist \vert_{\plane}(\closestlocalword{\plane}, \closestglobalword) \enspace.
\end{equation*}
Recalling that
\[
  \rho(M) = \frac{1}{\blocklength m} \sum_{\plane} \reldist \vert_{\plane}(\testerinput, \closestlocalword{\plane})
\]
we get that
\begin{equation*}
\reldist(\testerinput) \leq \robust(\testerinput) + \frac{1}{\blocklength m} \sum_{\plane} \reldist \vert_{\plane}(\closestlocalword{\plane}, \closestglobalword) \enspace.
\end{equation*}
Since $L$ is a vertex cover, the set $\overline{L} = V \setminus L$ is an independent set. Since $\abs{L} < (m-1)\codedist$,
\[
  \abs{\overline{L}} > \blocklength m - (m-1)\codedist = (m-1)(\blocklength-\codedist) + \blocklength\enspace.
\]
By \expref{Lemma}{lemma:interp}, $\exists f^* \in C^m$ such that
$f^* \vert_{\plane} = \closestlocalword{\plane}$ for every
$\plane \in \overline{L}$. Thus,
\begin{equation*}
\reldist(\testerinput) \leq \robust(\testerinput) + \frac{1}{\blocklength m} \sum_{\plane} \reldist \vert_{\plane}(\closestlocalword{\plane}, f^*) = \robust(\testerinput) + \frac{1}{\blocklength m} \sum_{\plane \in L} \reldist \vert_{\plane}(\closestlocalword{\plane}, f^*) \leq \robust(\testerinput) + \frac{\abs{L}}{\blocklength m} \enspace.
\end{equation*}
By \expref{Lemma}{lemma:robust},
\[
  \robust(\testerinput) \geq \frac{(m-2)}{4 (m-1)} \frac{\codedist^{m-1}}{\blocklength^{m-1}} \frac{\abs{L}}{\blocklength m}\enspace.
\]
Therefore,
\[
  \frac{\abs{L}}{\blocklength m} \leq \frac{4(m-1)\blocklength^{m-1}}{(m-2)\codedist^{m-1}} \robust(\testerinput)
\]
and so
\begin{equation*}
\reldist(\testerinput) \leq \robust(\testerinput) + \frac{\abs{L}}{\blocklength m} \leq \robust(\testerinput) \left(1 +  \frac{4(m-1)\blocklength^{m-1}}{(m-2)\codedist^{m-1}}  \right) \implies \robust(\testerinput) \geq \frac{1}{1 +  \frac{4(m-1)\blocklength^{m-1}}{(m-2)\codedist^{m-1}}} \reldist(\testerinput) \enspace.
\end{equation*}
Thus, $\forall \testerinput$, $\robust(\testerinput) \geq \alpha \reldist(\testerinput)$, for
\[
  \alpha = \min \left (\frac{1}{1 +  \frac{4(m-1)\blocklength^{m-1}}{(m-2)\codedist^{m-1}}}, \frac{m-2}{4m} \frac{\codedist^m}{\blocklength^m}\right)\enspace.
\]
Since $m \geq 3$, we have that
\[
  \frac{1}{1 +  \frac{4(m-1)\blocklength^{m-1}}{(m-2)\codedist^{m-1}}} \geq \frac{1}{1 + 8\frac{\blocklength^{m-1}}{\codedist^{m-1}} } \geq \frac{\codedist^{m-1}}{9\blocklength^{m-1}}
\qquad
\text{and}
\qquad
  \frac{m-2}{4m} \frac{\codedist^m}{\blocklength^m}  \geq \frac{1}{12} \frac{\codedist^m}{\blocklength^m}\enspace.
\]
Therefore,
\[
  \alpha \geq \min\left(\frac{\codedist^{m-1}}{9\blocklength^{m-1}}, \frac{1}{12} \frac{\codedist^m}{\blocklength^m}\right) = \frac{1}{12} \frac{\codedist^m}{\blocklength^m}\enspace.\qedhere
  \]
  \end{proof}






%
%


%
%
%
%
%
%

%
%
%
\section{Other results}
\label{sec:other}

%
Below we prove several results that are incomparable to \expref{Theorem}{thm:pat-robustness}.       %

We have already shown in \expref{Theorem}{thm:pat-robustness} that $\tester$ is robust for $\alpha \geq \frac{1}{12} \left(\frac{\codedist}{\blocklength}\right)^m$. Most of the proof was dedicated to analyzing the test when the
%
number of  %
%
%
vertices of large degree, %
$L$, was less than $(m-1)\codedist$. In this same regime, we can prove an incomparable value for $\alpha$. Specifically, we can show that for every $\testerinput$ such that $\abs{L} < (m-1) \codedist$ it holds that
\[
  \robust(\testerinput) \geq \frac{1}{m + c} \cdot \reldist(\testerinput)\enspace,
\]
where $c$ is a constant.

\begin{claim}
\label{claim:m}
If $\abs{L} < (m-1)\codedist$, then $\robust(\testerinput) \geq \frac{1}{m + c} \cdot \reldist(\testerinput)$, for $c = 32/9$. Combining with \expref{Theorem}{thm:pat-robustness}, this implies that
\[
  \robust(\testerinput) \geq \max \left(\frac{1}{m + c},\  \frac{1}{12} \left(\frac{\codedist}{\blocklength}\right)^m\right) \cdot \reldist(\testerinput)
\]
when $\abs{L} < (m-1)\codedist$.
\end{claim}

\begin{proof}
Let $I$ be the set of planes that are not in $L$.
By the assumption $\abs{L} < (m-1)\codedist$, we have $\abs{I} > (m-1)(\blocklength-\codedist) + \blocklength$,
and thus, by \expref{Lemma}{lemma:interp} there exists
$\closestglobalword \in C^m$ such that $\closestglobalword \vert_{\plane} = \closestlocalword{\plane}$ for all $\plane \in I$.

Let $K = \{p \in \Field^{\blocklength} : \forall \plane \in I, p \notin \plane\}$ be the set of points that are not contained in any plane in $I$.
For each $b \in [m]$ let $I_b = \{i \in [\blocklength] : (b,i) \in I\}$.
It is clear that $\abs{I} = \sum_{b \in [m]} \abs{I_b}$,
and we can rewrite $K$ as $K = \{p \in \Field^{\blocklength} : p_b \notin I_b \ \forall b \in [m]\}$.
Therefore,
\begin{equation*}
\abs{K} = \prod_{b = 1}^m (\blocklength - \abs{I_b}) \leq \left(\blocklength - \frac{1}{m} \sum_{b = 1}^m \abs{I_b} \right)^m = \blocklength^m \left(1 - \frac{1}{\blocklength m} \sum_{b = 1}^m \abs{I_b} \right)^m = \blocklength^m \left(\frac{\abs{L}}{\blocklength m}\right)^m \enspace,
\end{equation*}
where the inequality uses the AM-GM inequality.
Now, we show that $\reldist(\testerinput,f) \leq (m+c)\cdot \robust(\testerinput)$.
We start by writing $\reldist(\testerinput,\closestglobalword)$ as follows.
\begin{equation*}
  \reldist(\testerinput,\closestglobalword)
    = \frac{1}{\blocklength^m}  \abs{ \{p : \testerinput(p) \neq \closestglobalword(p) \}}
    = \frac{1}{\blocklength^m}  \abs{ \{p \in K : \testerinput(p) \neq \closestglobalword(p) \}}
        + \frac{1}{\blocklength^m}  \abs{ \{p \notin K: \testerinput(p) \neq \closestglobalword(p) \}} \enspace.
\end{equation*}
The first term is upper bounded by ${\abs{K}}/{\blocklength^m}$, and so it is at most $\left({\abs{L}}/{(\blocklength m)}\right)^m$.
In order to bound the second term, note that for all $p \notin K$ there exists a plane $\plane_p \in I$ such that $p \in H_p$,
and thus, $\closestglobalword(p) = \closestlocalword{\plane_p}(p)$. Therefore,
\begin{eqnarray*}
  \frac{1}{\blocklength^m}  \abs{ \{p \notin K: \testerinput(p) \neq \closestglobalword(p) \}}
  & = & \frac{1}{\blocklength^m}  \abs{ \{p \notin K: \testerinput(p) \neq \closestlocalword{\plane_p}(p) \}} \\
  & \leq & \frac{1}{\blocklength^m}  \abs{ \{p \in [\blocklength]^m: \testerinput(p) \neq \closestlocalword{\plane_p}(p) \}} \\
  & \leq & \frac{1}{\blocklength^m}  \sum_{p \in [\blocklength]^m} \abs{\{ \plane : p \in H, \testerinput(p) \neq \closestlocalword{\plane}(p) \}} \\
  & = & m \cdot \robust(\testerinput) \enspace.
\end{eqnarray*}
This implies that
\begin{equation*}
  \reldist(\testerinput,\closestglobalword)
  \leq  \left(\frac{\abs{L}}{\blocklength m}\right)^m + m \cdot \robust(\testerinput)\enspace.
\end{equation*}
Next, using the bound $\abs{L} < (m-1)\codedist$ in the assumption of the claim,
as well as the bound
\[
  \frac{\abs{L}}{\blocklength m} \leq \robust(\testerinput) \cdot \frac{4(m-1)}{m-2} \cdot \frac{\blocklength^{m-1}}{\codedist^{m-1}}
\]
from \expref{Lemma}{lemma:robust}, we get that
\begin{eqnarray*}
  \reldist(\testerinput,\closestglobalword)
  & \leq & \left(\frac{(m-1)\codedist}{\blocklength m}\right)^{m-1} \cdot
  \left( \robust(\testerinput) \cdot \frac{4(m-1)}{m-2} \cdot \frac{\blocklength^{m-1}}{\codedist^{m-1}} \right) + m \cdot \robust(\testerinput) \\
  & = & \left( \left(1  - \frac{1}{m}\right)^{m} \cdot \frac{4m}{m-2} + m \right) \cdot \robust(\testerinput) \enspace.
\end{eqnarray*}
For $m \geq 3$ we get that $\reldist(\testerinput) \leq (m + 32/9) \robust(\testerinput)$, as required.
\end{proof}

\begin{remark}
In fact, by a slightly modified argument (writing $\robust(\testerinput)$ as the sum over the intersections of $k$ planes) we can prove that for $\abs{L} < (m-1)\codedist$
it holds that
\[
  \reldist(\testerinput) \leq \robust(\testerinput) \left( k + c_k \frac{\blocklength^{m-k}}{\codedist^{m-k}}\right)\enspace,
\]
where $c_k$ is a constant for a fixed $k \in [m]$. The proof of \expref{Theorem}{thm:pat-robustness} used $k = 1$.
\end{remark}

We can also show that when $\abs{L} < (m-1) \codedist$, we get a robustness of $\alpha = 1$ plus an additive term of $\codedist/\blocklength$. Note that $\codedist$ is the distance of the code, so when $\codedist = O(\blocklength)$, the additive term is not small.

\begin{claim}
If $\abs{L} < (m-1) \codedist$, then $\reldist(\testerinput) \leq \robust(\testerinput ) + \codedist/\blocklength$.
\end{claim}
\begin{proof}
In the proof of \expref{Theorem}{thm:pat-robustness}, we showed that if $\abs{L} < (m-1)\codedist$, then
\begin{equation*}
\reldist(\testerinput) \leq \robust(\testerinput) + \frac{\abs{L}}{\blocklength m} \leq  \robust(\testerinput) + \frac{(m-1)\codedist}{\blocklength m} \leq \robust(\testerinput) + \frac{\codedist}{\blocklength} \enspace. \qedhere
\end{equation*}
\end{proof}

Next, we observe that for any linear $C$ of dimension $k$ we have a similar claim without the constraint on $\abs{L}$, if we replace $\codedist/\blocklength$ with $(\blocklength-k)/\blocklength$


\begin{claim}
\label{claim:robustness-of-RS-using-dim}
For any linear code $C \subseteq \Field^{\blocklength}$ of dimension $k$, it holds that $\reldist(\testerinput) \leq \robust(\testerinput) +  (\blocklength-k)/\blocklength$ unconditionally for all $\testerinput \in \Field^{\blocklength}$.
\end{claim}

\begin{proof}
Define $v_b = \sum_{\plane = (b,i)} \hamdist \vert_{\plane}(\testerinput, \closestlocalword{\plane})$, and without loss of generality assume that $v_1 \leq v_2 \leq \dots \leq v_m$. Observe that
\begin{flalign*}
\robust(\testerinput) = \frac{1}{\blocklength^m m} \sum_{b = 1}^m v_b
\enspace.
\end{flalign*}
Let $S \subseteq [\blocklength]$ be any subset of the coordinates that are pivotal for $C$, \ie, for any $k$ values given to the coordinates in $S$ there exists a unique codeword in $C$ that agrees with these values on $S$.\footnote{Note that $S$ exists since $C$ is a linear subspace of $\Field^{\blocklength}$.}
Let us identify $S$ with the corresponding $(1,i)$ planes. Since $S$ is a pivotal set for $C$, there exists $\closestglobalword$ in $C^m$ such that $\closestglobalword \vert_{\plane} = \closestlocalword{\plane}$ for every $\plane$ in $S$. Therefore,
\begin{align*}
&\reldist(\testerinput) \leq \reldist(\testerinput, \closestglobalword) = \frac{1}{\blocklength^m} \sum_{\plane = (1,i)} \hamdist \vert_{\plane}(\testerinput, \closestglobalword) \leq \frac{1}{\blocklength^m} \sum_{\plane \in S} \hamdist \vert_{\plane}(\testerinput, \closestglobalword) +  \frac{1}{\blocklength^m} \left(\blocklength - \abs{S}\right)\blocklength^{m-1} \\
& = \frac{1}{\blocklength^m} \sum_{\plane \in S} \hamdist \vert_{\plane}(\testerinput, \closestlocalword{\plane}) +  \frac{\blocklength-k}{\blocklength} \leq \frac{1}{\blocklength^m} v_1 +  \frac{\blocklength-k}{\blocklength} \\ %
&\leq \frac{1}{\blocklength^m m}\sum_{b = 1}^m v_b + \frac{\blocklength-k}{\blocklength}  = \robust(\testerinput) +  \frac{\blocklength-k}{\blocklength} \enspace. \qedhere
\end{align*}
\end{proof}

In particular, if $C$ is an MDS code then $k = \blocklength - \codedist + 1$, and so we get the following corollary.

\begin{corollary}
\label{cor:robustness-of-RS}
If $C$ is an MDS code, then $\reldist(\testerinput) \leq \robust(\testerinput) +  (\codedist-1)/\blocklength$ unconditionally for all $\testerinput \in \Field^{\blocklength}$.
\end{corollary}


%
%
%
\section{Robustness vs.\ soundness}
\label{sec:soundness}

In \expref{Theorem}{thm:pat-robustness}, we study the robustness of the test $\tester$. We now compare the robustness of $\tester$ to the soundness of $\tester$, and show a somewhat tight result for the soundness of $\tester$.

Robustness is a stronger guarantee than soundness, as $\Pr[\tester^{\testerinput} \text{ rejects}] \geq \robust(\testerinput)$. Thus, if we show that $\robust(\testerinput) \geq \alpha \reldist(\testerinput)$, then we can deduce that $\Pr[\tester^{\testerinput}\text{ rejects}] \geq \alpha \reldist(\testerinput)$, which upper bounds the soundness error. However, the converse is not true.

In fact, unlike for robustness, we can prove a somewhat tight result on the soundness of $\tester$, stated below. Note that the query complexity of $\tester$ is $\blocklength^{m-1}$, which is $N^{1 - \frac{1}{m}}$ where $N \DefineEqual \blocklength^m$ is the block length of $C^{m}$, the code being tested.

%

\begin{claim}
\label{claim:accprob}
Let $\epsilon_0 = ((\blocklength - \codedist)(m-1) + \blocklength )/\blocklength m$. If $\Pr[\tester^{\testerinput} \text{ accepts}] = \epsilon > \epsilon_0$, then $\reldist(\testerinput) \leq (1 - \epsilon)^m$.
\end{claim}
\begin{proof}
  Observe that
  \[
    \Pr[\tester^{\testerinput} \text{ accepts}] = \epsilon = \frac{1}{\blocklength m} \cdot \abs{I}\enspace,
  \]
  where $I$ is the set of planes $\plane$ such that $\testerinput \vert_{\plane} \in C^{m-1}$. The condition $\epsilon > \epsilon_0$ implies that $\abs{I} > (\blocklength - \codedist)(m-1) + \blocklength$, so by \expref{Lemma}{lemma:interp}, there exists an $\closestglobalword$ in $C^m$ such that $\closestglobalword \vert_{\plane} = \testerinput \vert_{\plane}$. Let $K = \{p : p \notin \plane \ \forall \plane \in I\}$. By the same logic as in the proof of \expref{Claim}{claim:m},
\begin{equation*}
\reldist(\testerinput) \leq \reldist(\testerinput, f) \leq \frac{\abs{K}}{\blocklength^m} \leq \left(1 - \frac{\abs{I}}{\blocklength m} \right)^m = (1 - \epsilon)^m \enspace. \qedhere
\end{equation*}
\end{proof}

We remark that the above claim is somewhat tight, as there are functions where ${\Pr[\tester^{\testerinput} \text{ accepts}] = \epsilon}$ and $\reldist(\testerinput) = \Omega((1- \epsilon)^m)$, and furthermore some threshold is necessary, as we now explain.
\begin{itemize}

  \item \emph{Necessity of $\epsilon$.}
  Fix $\epsilon$, and let $S$ be a subset of $[\blocklength]$ of size $(1-\epsilon) n$. Let $\testerinput$ be a word in $[\blocklength]^m$ obtained from some codeword of $C^m$ by corrupting all values on the set $S^{m}$ to be uniformly random. It is clear that with high probability $\reldist(\testerinput) = \Omega(\abs{S}^m/\blocklength^m) = \Omega((1- \epsilon)^m)$, and that the test $\tester$ will accept if it reads the plane $(b, i)$ for any $b \in [m]$ and $i \notin S$. Furthermore, the test $\tester$ will reject when it reads a plane $(b,i)$ for any $b \in [m]$ and $i \in S$. Therefore, $\Pr[\tester^{\testerinput} \text{ accepts}] = 1 - \abs{S}/\blocklength = \epsilon$.

  \item \emph{Necessity of $\epsilon_0$.}
Some threshold is necessary, as we can get an acceptance probability of $1/m$ ``for free'' by setting $\testerinput |_{(1,i)}$ to be a random $C^{m-1}$ codeword. Then, the test $\tester$ will accept with probability $1$ when it reads a $(1,i)$ plane, and will reject with high probability when it reads a $(b,i)$ plane for $b \ne 1$. However, $\testerinput$ will be very far from $C^m$, so having an acceptance probability of $1/m$ says little about $\reldist(\testerinput)$.
\end{itemize}



%
%
%
\section*{Acknowledgements}

The authors thank Eli Ben-Sasson, Oded Goldreich, and Madhu Sudan for helpful discussions, and the anonymous referees for their very useful comments. This work was supported in part by the Center for Long-Term Cybersecurity at UC Berkeley.

\nocite{conf-version}

\bibliographystyle{tocplain}   %

%
%
%
%
\bibliography{v016a005}

%
%
%
%
%
\begin{tocauthors}
\begin{tocinfo}[chiesa]
  Alessandro Chiesa\\
  Assistant professor\\
  Department of Electrical Engineering and Computer Science\\
  UC Berkeley\\
  Berkeley, CA, USA\\
  alexch\tocat{}berkeley\tocdot{}edu \\   %
  \url{https://people.eecs.berkeley.edu/~alexch/}      %
\end{tocinfo}
\begin{tocinfo}[manohar]
  Peter Manohar\\
  \phd\ student \\
  Computer Science Department\\
  Carnegie Mellon University\\
  Pittsburgh, PA, USA\\
  pmanohar\tocat{}cs\tocdot{}cmu\tocdot{}edu\\
  \url{https://www.cs.cmu.edu/~pmanohar}
\end{tocinfo}
\begin{tocinfo}[shinkar]
  Igor Shinkar\\
  Assistant professor\\
  School of Computing Science\\
  Simon Fraser University\\
  Burnaby, B.C. Canada\\
  ishinkar\tocat{}sfu\tocdot{}ca \\   %
  \url{https://www.cs.sfu.ca/~ishinkar/}
\end{tocinfo}
\end{tocauthors}

%
%
%
%
%
%
%
%
%
%
%
%
%
%
\begin{tocaboutauthors}
\begin{tocabout}[chiesa]  %
  \textsc{Alessandro Chiesa}
  graduated from MIT in 2014; his advisor was Silvio Micali.
  He was a postdoc at ETH Zurich between 2014-2015. Since then, he has been a
  faculty member at the University of California, Berkeley. He is interested
  in all kinds of proof systems, information theoretic and cryptographic.
\end{tocabout}
\begin{tocabout}[manohar]  %
  \textsc{Peter Manohar}
  is a \phd\ student at Carnegie Mellon University, advised by
  Venkat Guruswami. He graduated from UC Berkeley with a B.\,S. in EECS in
  2019 where he was advised by Professors Alessandro Chiesa and Ren Ng.
  He is broadly interested in complexity theory and cryptography, specifically
  in topics such as probabilistically checkable proofs, property testing,
  coding theory, and sum-of-squares.
\end{tocabout}
\begin{tocabout}[shinkar]
  \textsc{Igor Shinkar} is an assistant professor in the School of
Computing Science %
at   %
Simon Fraser University, Canada. He obtained his \phd\ in
Computer Science from
the   %
Weizmann Institute of Science, Israel in 2014 under the supervision
of %
%
Irit Dinur. He spent two years as a postdoc
%
at    %
the Courant Institute of Mathematical Sciences, NYU, and two years as
a postdoc at UC Berkeley. His research is in theoretical computer science,
discrete mathematics, probability, and the interplay between them.
\end{tocabout}
\end{tocaboutauthors}

\end{document}
